% =============================================================================
% Hauptdatei: Hausarbeit Experimentelle Röntgenphysik
% Kompilieren mit: pdflatex + biber (oder auf Overleaf: einfach hochladen)
% =============================================================================
\documentclass[
  12pt,
  a4paper,
  twoside=false,
  bibliography=totoc,
  listof=totoc,
  numbers=noenddot
]{scrreprt}

% ── Sprache und Encoding ────────────────────────────────────────────────────
\usepackage[utf8]{inputenc}
\usepackage[T1]{fontenc}
\usepackage[ngerman]{babel}

% ── Schrift ─────────────────────────────────────────────────────────────────
\usepackage{lmodern}
\usepackage{microtype}

% ── Seitenränder ────────────────────────────────────────────────────────────
\usepackage[
  left=2.5cm,
  right=2.5cm,
  top=2.5cm,
  bottom=2.5cm
]{geometry}

% ── Mathematik ──────────────────────────────────────────────────────────────
\usepackage{amsmath,amssymb,amsfonts}
\usepackage{bm}

% ── Grafiken und Farbe ─────────────────────────────────────────────────────
\usepackage{graphicx}
\graphicspath{{figures/}}
\usepackage{xcolor}
\usepackage{tikz}
\usetikzlibrary{arrows.meta, positioning, shapes.geometric, calc, decorations.pathmorphing}

% ── Tabellen ────────────────────────────────────────────────────────────────
\usepackage{booktabs}
\usepackage{tabularx}
\usepackage{multirow}
\usepackage{siunitx}
\sisetup{
  locale = DE,
  per-mode = symbol,
  separate-uncertainty = true
}

% ── Floats ──────────────────────────────────────────────────────────────────
\usepackage{float}
\usepackage[font=small, labelfont=bf, format=plain]{caption}
\usepackage{subcaption}

% ── Literatur ───────────────────────────────────────────────────────────────
\usepackage[
  backend=biber,
  style=numeric-comp,
  sorting=none,
  maxbibnames=6,
  minbibnames=3,
  doi=true,
  isbn=false,
  url=false
]{biblatex}
\addbibresource{references.bib}

% ── Hyperlinks ──────────────────────────────────────────────────────────────
\usepackage[
  colorlinks=true,
  linkcolor=black,
  citecolor=blue!60!black,
  urlcolor=blue!70!black,
  pdfauthor={},
  pdftitle={Strukturanalyse mit EvAX},
  bookmarksopen=true
]{hyperref}

% ── Kopf-/Fußzeilen ────────────────────────────────────────────────────────
\usepackage[headsepline]{scrlayer-scrpage}
\pagestyle{scrheadings}
\clearpairofpagestyles
\ohead{\pagemark}
\ihead{\headmark}

% ── Sonstiges ───────────────────────────────────────────────────────────────
\usepackage{enumitem}
\usepackage{csquotes}
\usepackage{listings}
\lstset{
  basicstyle=\ttfamily\small,
  frame=single,
  breaklines=true,
  captionpos=b,
  numbers=left,
  numberstyle=\tiny\color{gray},
  keywordstyle=\color{blue!70!black},
  commentstyle=\color{green!50!black},
}

% ── Eigene Befehle ──────────────────────────────────────────────────────────
\newcommand{\evax}{\textsc{EvAX}}
\newcommand{\feff}{\textsc{FEFF}}
\newcommand{\chik}{\ensuremath{\chi(k)}}
\newcommand{\chir}{\ensuremath{\chi(R)}}
\newcommand{\kweight}[1]{\ensuremath{k^{#1}}}
\newcommand{\Angstrom}{\text{\AA}}
\newcommand{\TODO}[1]{\textcolor{red}{\textbf{[TODO: #1]}}}
\newcommand{\PLACEHOLDER}[1]{\textcolor{blue!70!black}{\textit{#1}}}

% Kapiteltiefe im Inhaltsverzeichnis
\setcounter{tocdepth}{2}
\setcounter{secnumdepth}{3}

% ── Absatzformatierung ──────────────────────────────────────────────────────
\setlength{\parindent}{0pt}
\setlength{\parskip}{0.5\baselineskip}

% =============================================================================
\begin{document}

% ── Deckblatt ───────────────────────────────────────────────────────────────
% =============================================================================
% Deckblatt
% =============================================================================
\begin{titlepage}
  \centering

  \vspace*{1cm}

  {\large
    \TODO{Bergische Universität Wuppertal}\\[0.3em]
    \TODO{Fakultät für Mathematik und Naturwissenschaften}\\[0.3em]
    Fachgruppe Physik
  }

  \vspace{1.5cm}

  \rule{\textwidth}{0.8pt}

  \vspace{0.8cm}

  {\LARGE\bfseries
    Lokale Strukturanalyse von Legierungen\\
    und Nanopartikeln mittels EXAFS\\
    und Reverse-Monte-Carlo-Modellierung
  }

  \vspace{0.5cm}

  {\large
    Auswertung mit der Software \evax{}
  }

  \vspace{0.6cm}

  \rule{\textwidth}{0.8pt}

  \vspace{1.5cm}

  {\large Hausarbeit zur Vorlesung}\\[0.4em]
  {\Large\bfseries Experimentelle Röntgenphysik}

  \vspace{2cm}

  \begin{tabular}{r@{\hspace{1em}}l@{\hspace{3em}}r@{\hspace{1em}}l}
    \textbf{Autor 1:} & \TODO{Vorname Nachname} &
    \textbf{Matr.-Nr.:} & \TODO{000\,000} \\[0.6em]
    \textbf{Autor 2:} & \TODO{Vorname Nachname} &
    \textbf{Matr.-Nr.:} & \TODO{000\,000} \\[0.6em]
    \textbf{Autor 3:} & \TODO{Vorname Nachname} &
    \textbf{Matr.-Nr.:} & \TODO{000\,000} \\
  \end{tabular}

  \vspace{2cm}

  \begin{tabular}{rl}
    \textbf{Betreuer:}      & \TODO{Prof.\ Dr.\ Vorname Nachname} \\[0.4em]
    \textbf{Abgabedatum:}   & \TODO{TT.\ Monat JJJJ} \\
  \end{tabular}

  \vfill

  {\small Sommersemester 2026}

\end{titlepage}
\clearpage


% ── Abstract ────────────────────────────────────────────────────────────────
% =============================================================================
% Abstract
% =============================================================================
\chapter*{Zusammenfassung}
\addcontentsline{toc}{chapter}{Zusammenfassung}

Die Röntgenabsorptionsspektroskopie (XAS), insbesondere die Analyse der
feinen Struktur oberhalb der Absorptionskante (EXAFS), liefert elementspezifische
Informationen über die lokale atomare Umgebung in kondensierten Materialien.
Bei Systemen mit mehreren chemischen Komponenten -- etwa ternären Legierungen
oder bimetallischen Nanopartikeln -- stößt die konventionelle Kurvenanpassung
jedoch an ihre Grenzen, da die Zahl der zu bestimmenden Strukturparameter
rasch die durch die Daten verfügbare Information übersteigt.

Die vorliegende Arbeit setzt die Software \evax{} ein, die auf einem
evolutionären Algorithmus in Kombination mit Reverse Monte Carlo (RMC)
basiert, um dieses Problem zu umgehen.
\evax{} erzeugt dreidimensionale Strukturmodelle und optimiert die atomaren
Positionen so, dass die simulierten EXAFS-Spektren aller beteiligten Kanten
simultan mit den experimentellen Daten übereinstimmen.

Im Mittelpunkt stehen zwei Proben:
eine \TODO{CoNiCu}-Mittelentropielegierung (MEA), deren drei K-Kanten
gleichzeitig ausgewertet werden, sowie \TODO{CoPt}-Nanopartikel, bei denen
die Verteilung der Elemente auf die Gitterplätze untersucht wird.
Für beide Systeme werden radiale Verteilungsfunktionen, Koordinationszahlen
und die chemische Nahordnung (Warren-Cowley-Parameter) bestimmt.
Darüber hinaus wird durch systematische Variation der Fit-Parameter und
den Vergleich unterschiedlicher Ausgangsstrukturen (FCC, BCC, HCP)
die Robustheit der Ergebnisse überprüft.

\vspace{0.5\baselineskip}

\textbf{Stichworte:}
EXAFS, Reverse Monte Carlo, evolutionärer Algorithmus, \evax{},
lokale Struktur, chemische Nahordnung, Mittelentropielegierung, Nanopartikel


% ── Inhaltsverzeichnis ──────────────────────────────────────────────────────
\tableofcontents
\clearpage

% ── Kapitel ─────────────────────────────────────────────────────────────────
% =============================================================================
% Kapitel 1: Einleitung
% =============================================================================
\chapter{Einleitung}
\label{ch:einleitung}

Die gezielte Entwicklung neuer Funktionsmaterialien setzt ein genaues
Verständnis ihrer atomaren Struktur voraus.
Während Beugungsmethoden wie Röntgendiffraktometrie die langreichweitige
Kristallstruktur beschreiben, fehlt ihnen in der Regel die Empfindlichkeit
für lokale Abweichungen von der mittleren Ordnung.
Gerade bei Mehrkomponentensystemen -- Hochentropie- und
Mittelentropielegierungen, bimetallischen Nanopartikeln oder dotierten
Oxiden -- bestimmen jedoch diese lokalen Arrangements wesentlich die
mechanischen, katalytischen und magnetischen Eigenschaften
\cite{Yeh2004, George2019, Ferrando2008}.

Die Röntgenabsorptionsfeinstrukturspektroskopie (EXAFS) bietet hier einen
komplementären Zugang.
Sie nutzt die Modulation des Absorptionskoeffizienten oberhalb einer
elementspezifischen Absorptionskante, um Bindungslängen,
Koordinationszahlen und die thermische bzw.\ statische Unordnung in der
unmittelbaren Nachbarschaft eines Absorberatoms zu bestimmen
\cite{Stern1974, Rehr2000}.
Da jede Kante einzeln adressiert wird, lässt sich die lokale Umgebung
jedes Elements in einem Mehrstoffsystem separat untersuchen.

\section{Motivation}

Bei einer ternären Legierung wie CoNiCu existieren neun verschiedene
Nachbarpaare (Co-Co, Co-Ni, Co-Cu, Ni-Co,\,\ldots).
Jeder dieser Beiträge wird durch mindestens vier Parameter beschrieben
(Koordinationszahl, Abstand, Debye-Waller-Faktor, Energieverschiebung),
sodass bereits für die erste Koordinationsschale $9 \times 4 = 36$
Fitparameter anfallen.
Die Menge an unabhängiger Information, die ein EXAFS-Spektrum enthält,
ist jedoch durch die sogenannte Nyquist-Zahl
\begin{equation}
  N_{\text{ind}} = \frac{2\,\Delta k\,\Delta R}{\pi} + 2
  \label{eq:nyquist}
\end{equation}
begrenzt, wobei $\Delta k$ den nutzbaren $k$-Bereich und $\Delta R$ das
betrachtete $R$-Fenster bezeichnen \cite{Stern1993}.
Für typische Werte ($\Delta k \approx \SI{8}{\per\Angstrom}$,
$\Delta R \approx \SI{2}{\Angstrom}$) ergibt sich
$N_{\text{ind}} \approx 12$ pro Kante, also insgesamt rund 36 unabhängige
Datenpunkte für drei Kanten -- kaum mehr als die Zahl der Parameter.
Ein konventioneller Least-Squares-Fit wird unter diesen Bedingungen
numerisch instabil und physikalisch mehrdeutig.

\section{Der RMC/EA-Ansatz}

Eine alternative Strategie kehrt das Fitproblem um:
Statt analytische Funktionen an das Spektrum anzupassen, wird ein
atomistisches Strukturmodell konstruiert und solange verändert, bis die
daraus berechneten Spektren aller Kanten gleichzeitig mit dem Experiment
übereinstimmen.
Die Software \evax{} implementiert diesen Ansatz als Kombination eines
evolutionären Algorithmus (EA) mit Reverse Monte Carlo (RMC)
\cite{Timoshenko2009, Timoshenko2012a, Timoshenko2014}.
Das Verfahren kommt ohne explizite Parametrisierung der
Paarverteilungsfunktionen aus und liefert unmittelbar
dreidimensionale Strukturmodelle, aus denen sich Koordinationszahlen,
Bindungslängen und die chemische Nahordnung ableiten lassen.

\section{Zielsetzung}

Ziel dieser Arbeit ist es, \evax{} auf zwei experimentelle Datensätze
anzuwenden:

\begin{enumerate}[label=(\roman*)]
  \item Eine \textbf{\TODO{CoNiCu}-Mittelentropielegierung (MEA)}, gemessen
        an den K-Kanten von Cobalt, Nickel und Kupfer.
        Durch gleichzeitige Auswertung aller drei Kanten sollen die
        elementaufgelösten Koordinationszahlen und insbesondere die
        Warren-Cowley-Parameter der chemischen Nahordnung bestimmt werden.
  \item \textbf{\TODO{CoPt}-Nanopartikel}, bei denen die Frage im
        Vordergrund steht, ob die beiden Elemente statistisch auf
        einem gemeinsamen Gitter verteilt sind oder eine geordnete
        Phase (etwa L1$_0$) bilden.
\end{enumerate}

Darüber hinaus wird die Sensitivität der Ergebnisse gegenüber den
Fit-Parametern ($k$-Bereich, $k$-Gewichtung, Vergleichsraum) durch
eine systematische Parameterstudie untersucht, und die Wahl der
Kristallstruktur (FCC vs.\ BCC vs.\ HCP) wird durch Vergleich der
Fitresiduen validiert.

\section{Aufbau der Arbeit}

Kapitel~\ref{ch:theorie} legt die physikalischen Grundlagen der
EXAFS-Spektroskopie dar.
In Kapitel~\ref{ch:evax} wird die Software \evax{} vorgestellt und ihr
Arbeitsablauf erläutert.
Kapitel~\ref{ch:methodik} beschreibt die konkrete Vorgehensweise,
einschließlich der Datenvorverarbeitung und der gewählten Fit-Strategie.
Die Ergebnisse der Auswertung beider Proben werden in
Kapitel~\ref{ch:auswertung} zusammengefasst und in
Kapitel~\ref{ch:diskussion} diskutiert.
Kapitel~\ref{ch:fazit} schließt die Arbeit mit einem Fazit und Ausblick ab.

% =============================================================================
% Kapitel 2: Theoretische Grundlagen
% =============================================================================
\chapter{Theoretische Grundlagen}
\label{ch:theorie}

\section{Röntgenabsorption und die Absorptionskante}

Trifft ein monochromatischer Röntgenstrahl der Intensität $I_0$ auf eine
Probe der Dicke $d$, so wird die transmittierte Intensität durch das
Lambert-Beersche Gesetz beschrieben:
\begin{equation}
  I(E) = I_0\,\mathrm{e}^{-\mu(E)\,d}\,,
  \label{eq:lambert}
\end{equation}
wobei $\mu(E)$ der lineare Absorptionskoeffizient ist.
Wird die Photonenenergie $E$ über eine Ionisierungsschwelle eines in der
Probe enthaltenen Elements hinweg erhöht, steigt $\mu(E)$ sprunghaft
an -- man spricht von einer \emph{Absorptionskante}.
An der K-Kante beispielsweise wird ein 1s-Elektron in das Kontinuum
angeregt; die Kantenenergie $E_0$ ist elementspezifisch und liegt für
die 3d-Übergangsmetalle im Bereich von ca.\ \SIrange{5}{9}{\kilo\electronvolt}
\cite{Koningsberger1988}.

Die Absorptionskante ist nicht nur ein Schwellenwert, sondern enthält in
ihrer näheren Umgebung eine reiche Struktur:
Der Bereich bis etwa \SI{30}{\electronvolt} oberhalb von $E_0$ wird
als XANES (X-ray Absorption Near Edge Structure) bezeichnet und
spiegelt die elektronische Struktur und Symmetrie des Absorbers wider.
Der sich anschließende Bereich bis zu einigen hundert Elektronenvolt
oberhalb der Kante heißt EXAFS (Extended X-ray Absorption Fine
Structure) und wird durch die Interferenz des emittierten
Photoelektrons mit seinen an den Nachbaratomen gestreuten Anteilen
verursacht.

% XAS-Prinzipskizze (TikZ)
% Standalone TikZ figure: XAS measurement principle
% Can be compiled separately or included via \input
\begin{figure}[htbp]
  \centering
  \begin{tikzpicture}[
    scale=0.9,
    atom/.style={circle, draw, fill=blue!20, minimum size=0.5cm,
                 inner sep=0pt, font=\tiny},
    absorber/.style={circle, draw, fill=red!40, minimum size=0.6cm,
                     inner sep=0pt, font=\tiny\bfseries},
    photon/.style={-{Stealth}, decorate,
                   decoration={snake, amplitude=1.5pt, segment length=6pt},
                   thick, red!70!black},
    electron/.style={-{Stealth}, thick, blue!70!black, dashed},
    scattered/.style={-{Stealth}, thick, green!50!black, dashed},
  ]
    % X-ray beam
    \draw[photon] (-3, 0) -- (-0.5, 0)
      node[midway, above=4pt, font=\small, text=red!70!black] {$h\nu$};

    % Absorber atom
    \node[absorber] (abs) at (0,0) {A};

    % Neighbor atoms
    \node[atom] (n1) at (2.2, 0.8) {N};
    \node[atom] (n2) at (2.0, -1.0) {N};
    \node[atom] (n3) at (-1.0, 2.0) {N};
    \node[atom] (n4) at (1.5, 2.2) {N};

    % Photoelectron wave (outgoing)
    \draw[electron] (abs.east) -- (n1.west)
      node[midway, above, font=\scriptsize] {$e^-$};
    \draw[electron] (abs.east) -- (n2.north west);

    % Scattered wave (returning)
    \draw[scattered] (n1.south west) -- (abs.north east);
    \draw[scattered] (n2.north) -- (abs.south east);

    % Interference annotation
    \draw[<->, thick, gray] (0, -2) -- (2.2, -2)
      node[midway, below, font=\small] {$R_j$};

    % Labels
    \node[font=\small, text=red!70!black, below left] at (-0.3, -0.5)
      {Absorber};
    \node[font=\small, text=blue!70!black, right] at (2.5, 0.8)
      {Nachbar};

    % Right side: schematic spectrum
    \begin{scope}[xshift=7cm, yshift=-0.5cm]
      \draw[->, thick] (0, 0) -- (4.5, 0)
        node[right, font=\small] {$E$};
      \draw[->, thick] (0, 0) -- (0, 3.5)
        node[above, font=\small] {$\mu(E)$};

      % Pre-edge
      \draw[thick, blue!60!black] (0.2, 0.5) -- (1.2, 0.6);

      % Edge jump
      \draw[thick, blue!60!black] (1.2, 0.6) -- (1.4, 2.5);

      % EXAFS oscillations
      \draw[thick, blue!60!black]
        (1.4, 2.5) .. controls (1.7, 2.8) and (1.9, 2.2) ..
        (2.1, 2.5) .. controls (2.3, 2.8) and (2.5, 2.3) ..
        (2.7, 2.5) .. controls (2.9, 2.7) and (3.1, 2.4) ..
        (3.3, 2.5) .. controls (3.5, 2.6) and (3.7, 2.45) ..
        (4.0, 2.5);

      % Annotations
      \draw[<->, gray, thick] (1.15, 0.6) -- (1.15, 2.5)
        node[midway, left, font=\scriptsize] {$\Delta\mu_0$};
      \node[font=\scriptsize] at (1.3, 0.2) {$E_0$};
      \draw[gray, dashed] (1.3, 0) -- (1.3, 0.6);

      \draw[decorate, decoration={brace, amplitude=4pt}]
        (1.5, 3.2) -- (4.0, 3.2)
        node[midway, above=4pt, font=\scriptsize] {EXAFS};

      \draw[decorate, decoration={brace, amplitude=4pt}]
        (1.1, 3.2) -- (1.5, 3.2)
        node[midway, above=4pt, font=\scriptsize] {XANES};
    \end{scope}
  \end{tikzpicture}
  \caption{Prinzip der Röntgenabsorptionsspektroskopie.
    Links: Das einfallende Photon erzeugt am Absorberatom (A) ein
    Photoelektron, das an den Nachbaratomen (N) gestreut wird.
    Die Interferenz zwischen ausgehender und rückgestreuter Welle
    moduliert den Absorptionskoeffizienten.
    Rechts: Schematisches Absorptionsspektrum $\mu(E)$ mit
    Kantensprung $\Delta\mu_0$, XANES- und EXAFS-Bereich.}
  \label{fig:xas-principle}
\end{figure}


\section{Die EXAFS-Gleichung}
\label{sec:exafs-gleichung}

Die EXAFS-Oszillation wird als normierter, kantenbereinigte Anteil des
Absorptionskoeffizienten definiert:
\begin{equation}
  \chi(E) = \frac{\mu(E) - \mu_0(E)}{\Delta\mu_0(E_0)}\,,
  \label{eq:chi_E}
\end{equation}
mit dem atomaren Hintergrund $\mu_0(E)$ und dem Kantensprung
$\Delta\mu_0(E_0)$.
Die Umrechnung in den Wellenzahlraum erfolgt über
\begin{equation}
  k = \sqrt{\frac{2m_e(E - E_0)}{\hbar^2}}\,,
  \label{eq:k_def}
\end{equation}
sodass das EXAFS-Signal als Funktion von $k$ geschrieben werden kann.

Im Einfachstreumodell ergibt sich die bekannte EXAFS-Gleichung
\cite{Stern1974, Sayers1971}:
\begin{equation}
  \chi(k) = S_0^2 \sum_j \frac{N_j\,|f_j(k)|}{k\,R_j^2}\,
    \sin\!\bigl(2kR_j + \delta_j(k)\bigr)\,
    \mathrm{e}^{-2\sigma_j^2 k^2}\,
    \mathrm{e}^{-2R_j/\lambda(k)}\,.
  \label{eq:exafs}
\end{equation}

Die Summation läuft über alle Koordinationsschalen $j$ mit folgenden Größen:
\begin{itemize}[nosep]
  \item $N_j$ -- Koordinationszahl (Anzahl der Nachbarn in Schale $j$),
  \item $R_j$ -- mittlerer Abstand zum Absorber,
  \item $f_j(k)$ und $\delta_j(k)$ -- Rückstreuamplitude und
        Phasenverschiebung, die vom Streuerelement abhängen,
  \item $\sigma_j^2$ -- mittlere quadratische Abweichung des Abstands
        (Debye-Waller-Faktor), die sowohl thermische Schwingungen als auch
        statische Unordnung erfasst,
  \item $\lambda(k)$ -- mittlere freie Weglänge des Photoelektrons,
  \item $S_0^2$ -- Amplitudenreduktionsfaktor, der inelastische
        Vielteilcheneffekte berücksichtigt.
\end{itemize}

Diese Gleichung verdeutlicht den zentralen Informationsgehalt eines
EXAFS-Spektrums: Jede Nachbarschale erzeugt eine Sinusschwingung im
$k$-Raum, deren Frequenz vom Abstand $R_j$ und deren Amplitude von der
Koordinationszahl $N_j$, der Rückstreuamplitude und dem
Debye-Waller-Faktor abhängt.

\section{Fouriertransformation und $R$-Raum}

Durch Fouriertransformation von $\chi(k)$ in den $R$-Raum lassen sich
die Beiträge verschiedener Koordinationsschalen separieren.
Üblicherweise wird das $k$-gewichtete Signal $\chi(k)\cdot k^n$ mit
einer Fensterfunktion $W(k)$ multipliziert:
\begin{equation}
  \tilde{\chi}(R) = \frac{1}{\sqrt{2\pi}}
    \int_{k_{\min}}^{k_{\max}} \chi(k)\,k^n\,W(k)\,
    \mathrm{e}^{2ikR}\,\mathrm{d}k\,.
  \label{eq:ft}
\end{equation}
Der Betrag $|\tilde{\chi}(R)|$ zeigt Maxima bei Abständen, die in erster
Näherung den Positionen der Koordinationsschalen entsprechen.
Zu beachten ist, dass die Peaks gegenüber den tatsächlichen
Bindungslängen um etwa
$\SIrange{0.2}{0.5}{\Angstrom}$ zu kleineren Werten verschoben sind,
da die Phasenverschiebung $\delta_j(k)$ nicht berücksichtigt wird.

Die Wahl der $k$-Gewichtung beeinflusst die relative Empfindlichkeit:
$k^1$ betont leichte Rückstreuer bei kleinen $k$-Werten,
$k^3$ hebt schwere Streuer bei hohen $k$ hervor,
und $k^2$ stellt einen häufig gewählten Kompromiss dar \cite{Bunker2010}.

\section{Mehrfachstreuung}

Gleichung~\eqref{eq:exafs} beschreibt ausschließlich die Einfachstreuung
(single scattering, SS), bei der das Photoelektron auf direktem Weg zum
Nachbaratom und zurück streut.
Bei kollinearen oder nahezu kollinearen Atomanordnungen können jedoch
Mehrfachstreupfade (multiple scattering, MS) einen erheblichen Beitrag
leisten.
Insbesondere Drei- und Vierbeinpfade, bei denen das Elektron an
Zwischenatomen umgelenkt wird, sind in dicht gepackten Strukturen wie
FCC-Metallen von Bedeutung, da dort die Fokussierung entlang der
$\langle 110\rangle$-Richtung den Effekt verstärkt \cite{Zabinsky1995}.

Die Software \feff{} \cite{Rehr2010} berechnet sowohl Einfach- als auch
Mehrfachstreupfade ab initio unter Verwendung der
Muffin-Tin-Approximation für das Kristallpotential und einer
selbstkonsistenten Berechnung der Streuphasen.
Die Berücksichtigung von MS-Pfaden ist insbesondere bei der Analyse
höherer Koordinationsschalen unerlässlich.

\section{Wavelet-Transformation}
\label{sec:wavelet}

Eine konventionelle Fouriertransformation liefert nur $R$-aufgelöste
Information; die Zuordnung eines Peaks zu einem bestimmten Streuerelement
ist allein aus $|\tilde{\chi}(R)|$ nicht möglich, da die $k$-Abhängigkeit
verloren geht.
Die Wavelet-Transformation (WT) umgeht diese Einschränkung, indem sie
eine zweidimensionale Darstellung in $(k,R)$ erzeugt
\cite{Funke2005, Timoshenko2009}.
Dadurch lassen sich Streubeiträge bei gleichem Abstand, aber
unterschiedlichen Rückstreuelementen (und damit unterschiedlicher
$k$-Abhängigkeit von $f_j(k)$) voneinander trennen.

In \evax{} wird der Vergleich zwischen Theorie und Experiment wahlweise
im $k$-Raum, im $R$-Raum oder im Wavelet-Raum durchgeführt.
Der Wavelet-Raum stellt dabei die schärfste Bedingung an die
Übereinstimmung, da er beide Dimensionen gleichzeitig berücksichtigt
\cite{Timoshenko2014}.

\section{Chemische Nahordnung: der Warren-Cowley-Parameter}
\label{sec:sro}

In einer Legierung mit mehreren Komponenten interessiert nicht nur die
mittlere Struktur, sondern auch die Frage, ob bestimmte
Elementkombinationen als Nachbarn bevorzugt oder vermieden werden.
Der Warren-Cowley-Parameter quantifiziert diese chemische Nahordnung
(short-range order, SRO) \cite{Warren1969, Cowley1950}:
\begin{equation}
  \alpha_{ij} = 1 - \frac{P_{ij}}{c_j}\,,
  \label{eq:warren_cowley}
\end{equation}
wobei $P_{ij}$ die bedingte Wahrscheinlichkeit ist, in der nächsten
Nachbarschale eines $i$-Atoms ein $j$-Atom zu finden, und $c_j$ der
Molenbruch von Element $j$ ist.

Die Interpretation ist unmittelbar anschaulich:
\begin{itemize}[nosep]
  \item $\alpha_{ij} = 0$: statistisch zufällige Verteilung (ideale Mischkristall),
  \item $\alpha_{ij} > 0$: $j$-Nachbarn werden in der Umgebung von $i$ vermieden,
  \item $\alpha_{ij} < 0$: $j$-Nachbarn werden bevorzugt.
\end{itemize}

Da \evax{} ein vollständiges atomistisches Modell liefert, lassen sich
$P_{ij}$ und damit $\alpha_{ij}$ direkt aus der finalen Struktur
berechnen, ohne zusätzliche Annahmen über die Form der
Paarverteilungsfunktion.

\section{Grenzen der EXAFS-Methode}

Trotz ihrer Stärken unterliegt die EXAFS-Analyse einigen
grundsätzlichen Einschränkungen:
\begin{itemize}
  \item \textbf{Lokale Sonde:} EXAFS erfasst nur die gemittelte
        Nahordnung innerhalb weniger Ångström um das Absorberatom.
        Langreichweitige Ordnung (Domänengrößen, Gitterkonstanten im
        klassischen Sinn) ist der Methode nicht zugänglich.
  \item \textbf{Parameterkorrelation:} $N_j$ und $\sigma_j^2$ sind
        stark korreliert, was die gleichzeitige Bestimmung beider Größen
        aus einem einzelnen Spektrum erschwert \cite{Bunker2010}.
  \item \textbf{Begrenzte Informationsdichte:} Die Nyquist-Zahl
        (Gl.~\eqref{eq:nyquist}) limitiert die Zahl der zuverlässig
        bestimmbaren Parameter.
  \item \textbf{Mehrkomponentige Systeme:} Bei mehr als zwei
        Streuerelementen in der gleichen Schale werden die Beiträge
        zunehmend schwer trennbar, sofern nicht mehrere Kanten
        gleichzeitig ausgewertet werden.
\end{itemize}

Gerade der letzte Punkt motiviert den Einsatz von Methoden wie \evax{},
die durch simultane Multi-Edge-Auswertung die Informationsdichte
erheblich steigern.

% =============================================================================
% Kapitel 3: Die Software EvAX
% =============================================================================
\chapter{Die Software EvAX}
\label{ch:evax}

\section{Überblick und Einordnung}

\evax{} (Evolutionary Algorithm for XAS) ist ein Programmpaket zur
Analyse von EXAFS-Daten, das an der Universität Lettland entwickelt
wurde \cite{Timoshenko2009, Timoshenko2012a, Timoshenko2012b}.
Im Unterschied zu konventionellen Least-Squares-Verfahren
-- wie sie etwa in \textsc{Artemis}/\textsc{Demeter}
\cite{Ravel2005} implementiert sind -- arbeitet \evax{} nicht mit
einer parametrischen Beschreibung der radialen Verteilungsfunktionen,
sondern optimiert die atomaren Positionen in einem dreidimensionalen
Strukturmodell direkt.

Die zentrale Idee besteht darin, eine Superzelle mit einigen hundert
Atomen aufzubauen, für jeden Absorbertyp das zugehörige EXAFS-Signal
ab initio zu berechnen und die Übereinstimmung mit dem Experiment
durch iterative Verrückung einzelner Atome zu verbessern.
Die Kombination eines evolutionären Algorithmus mit Elementen der
Reverse-Monte-Carlo-Methode (RMC) stellt sicher, dass der
Konfigurationsraum effizient abgesucht wird und lokale Minima
überwunden werden können.

\section{Grundprinzip: RMC und evolutionärer Algorithmus}

\subsection{Reverse Monte Carlo (RMC)}

Die RMC-Methode \cite{McGreevy1988} geht von einem atomistischen
Startmodell aus und modifiziert es durch zufällige Verschiebungen
einzelner Atome.
Nach jeder Verrückung wird das resultierende Signal berechnet und
mit dem Experiment verglichen.
Verbessert sich die Übereinstimmung (gemessen durch ein
Residuenfunktional), wird die Verschiebung akzeptiert; verschlechtert
sie sich, wird die Verschiebung mit einer bestimmten
Wahrscheinlichkeit dennoch zugelassen -- analog zum
Metropolis-Kriterium in der statistischen Physik.
Dieser stochastische Akzeptanzmechanismus verhindert, dass der
Algorithmus in einem lokalen Minimum verharrt.

\subsection{Evolutionärer Algorithmus}

\evax{} erweitert das RMC-Prinzip um einen evolutionären Algorithmus:
Es werden nicht eine, sondern $N_{\text{states}}$ (typischerweise 16)
parallele Strukturkopien (\emph{States}) gleichzeitig optimiert.
In regelmäßigen Abständen werden die States anhand ihrer Fitgüte
bewertet; schlechte Strukturen werden durch Kopien guter Strukturen
ersetzt, wobei leichte Mutationen eingeführt werden.
Dieses Selektions-Mutations-Schema entspricht dem Prinzip der natürlichen
Auslese und erhöht die Wahrscheinlichkeit, das globale Minimum zu finden
\cite{Timoshenko2014}.

\section{FEFF-Integration und Spektrumberechnung}

Für die Berechnung des theoretischen EXAFS-Signals verwendet \evax{}
das Programm \feff{} (Version~8.5) \cite{Rehr2010} als externen
Rechenkern.
\feff{} bestimmt die Streuphasen und -amplituden ab initio aus der
Muffin-Tin-Approximation des Kristallpotentials und erzeugt die
Beiträge aller relevanten Einfach- und Mehrfachstreupfade.

Da eine vollständige \feff{}-Rechnung für alle 256~Atome bei jeder
Iteration viel zu aufwendig wäre, nutzt \evax{} eine
Clusterexpansion: Die lokale Umgebung jedes Absorberatoms wird mit
einer Bibliothek bereits berechneter Cluster verglichen, und nur bei
hinreichend starker Veränderung wird eine neue \feff{}-Rechnung
ausgelöst.
Der Parameter \texttt{Update\_basis\_every} steuert, nach wie vielen
Iterationen die Streubasis aktualisiert wird.

\section{Multi-Edge-Auswertung}

Ein besonderer Vorteil von \evax{} liegt in der simultanen Auswertung
mehrerer Absorptionskanten.
Bei einer ternären Legierung wie CoNiCu werden die K-Kanten aller drei
Elemente gleichzeitig gefittet.
Da alle Kanten vom selben Strukturmodell erzeugt werden, sind die
Fitparameter implizit verknüpft: Eine Verrückung eines Co-Atoms
beeinflusst sowohl das Co-K- als auch das Ni-K- und Cu-K-Spektrum.
Dies erhöht die effektive Informationsdichte erheblich und reduziert
die Mehrdeutigkeit des Fits \cite{Timoshenko2012a}.

Die relative Gewichtung der einzelnen Kanten im Residuenfunktional wird
durch \evax{} automatisch angepasst (\texttt{adjust\_spectra\_weights}),
sodass keine Kante den Fit dominiert.

\section{Wichtige Eingabeparameter}
\label{sec:evax-params}

Tabelle~\ref{tab:evax-params} fasst die zentralen Parameter zusammen,
die in der Eingabedatei (\texttt{parameters.dat}) spezifiziert werden.

\begin{table}[htbp]
  \centering
  \caption{Zentrale Eingabeparameter von \evax{} und ihre Bedeutung.}
  \label{tab:evax-params}
  \small
  \begin{tabularx}{\textwidth}{lX}
    \toprule
    \textbf{Parameter} & \textbf{Bedeutung} \\
    \midrule
    \texttt{File\_with\_structure}
      & Pfad zur Startstruktur im POSCAR-Format (Gittervektoren und
        Atompositionen) \\
    \texttt{Supercell}
      & Multiplikation der Einheitszelle, z.\,B.\ \texttt{1x1x1} für
        eine bereits aufgebaute Superzelle \\
    \texttt{S02}
      & Amplitudenreduktionsfaktor $S_0^2$ für jede Kante.
        Wird mit der \texttt{-autoES}-Option automatisch bestimmt \\
    \texttt{dE0}
      & Energieverschiebung $\Delta E_0$ (\si{\electronvolt}) zwischen
        theoretischer und experimenteller Kantenlage \\
    \texttt{k\_min}, \texttt{k\_max}
      & Grenzen des genutzten $k$-Bereichs (\si{\per\Angstrom}) \\
    \texttt{k\_power}
      & Exponent $n$ der $k$-Gewichtung ($k^n \cdot \chi(k)$) \\
    \texttt{Space}
      & Vergleichsraum: \texttt{k} ($k$-Raum), \texttt{R} ($R$-Raum)
        oder \texttt{w} (Wavelet) \\
    \texttt{N\_legs}
      & Maximale Zahl der Streuarme pro Pfad (2\,=\,nur
        Einfachstreuung, 4\,=\,inkl.\ MS) \\
    \texttt{R\_max\_for\_FEFF}
      & Maximaler Radius für die \feff{}-Cluster (\si{\Angstrom}) \\
    \texttt{Maximal\_step\_length}
      & Maximale Verschiebung pro Atom und Iteration (\si{\Angstrom}) \\
    \texttt{Maximal\_displacement}
      & Maximale kumulative Verschiebung eines Atoms gegenüber der
        Startposition (\si{\Angstrom}) \\
    \texttt{Number\_of\_states}
      & Anzahl paralleler Strukturkopien im EA \\
    \texttt{Stop\_after}
      & Zahl der Iterationen in der Screening-Phase \\
    \texttt{Froze\_in}
      & Abbruchkriterium: Iterationen ohne Verbesserung bis zur
        Konvergenz \\
    \bottomrule
  \end{tabularx}
\end{table}

\section{Ausgaben und Ergebnisgrößen}

Nach Abschluss der Optimierung liefert \evax{} folgende Ergebnisse:

\begin{description}[style=nextline]
  \item[Finale Struktur (\texttt{final.xyz})]
    Kartesische Koordinaten aller Atome in der optimierten Superzelle,
    aus denen Paarverteilungsfunktionen, Koordinationszahlen und SRO-Parameter
    berechnet werden können.
  \item[EXAFS-Spektren (\texttt{EXAFS.dat})]
    Berechnetes und experimentelles $\chi(k)$ für jede Kante, zur
    visuellen und quantitativen Beurteilung der Fitgüte.
  \item[Fouriertransformierte (\texttt{expFT.dat}, \texttt{expBFT.dat})]
    FT des experimentellen und berechneten Signals.
  \item[Radiale Verteilungsfunktionen (RDF)]
    Elementaufgelöste $g_{ij}(R)$ für alle Paarkombinationen.
  \item[Konvergenzverlauf (\texttt{output.dat})]
    Residuen als Funktion der Iterationszahl für jede Kante und als
    gewichtete Summe.
\end{description}

\section{Beispielhafter Workflow in EvAX}
\label{sec:workflow}

Im Folgenden wird der typische Arbeitsablauf bei einer \evax{}-Analyse
Schritt für Schritt beschrieben.
Abbildung~\ref{fig:workflow} zeigt eine schematische Darstellung.

\subsection{Datenvorverarbeitung}

Die experimentellen Rohdaten liegen als Absorptionsspektren $\mu(E)$
vor und müssen zunächst in $\chi(k)$-Daten umgewandelt werden.
Dies geschieht in der Regel mit der Software \textsc{Athena}
\cite{Ravel2005} oder mit der Python-Bibliothek \textsc{larch}
\cite{Newville2013} und umfasst:
\begin{enumerate}[nosep]
  \item Bestimmung der Kantenenergie $E_0$,
  \item Abzug des Vor-Kanten-Hintergrunds (Pre-Edge-Subtraktion),
  \item Normierung auf den Kantensprung $\Delta\mu_0$,
  \item Entfernung des atomaren Hintergrunds mittels Spline-Fit
        (\texttt{autobk}-Algorithmus),
  \item Export der $\chi(k)$-Daten.
\end{enumerate}

Bei Multi-Edge-Messungen, bei denen mehrere Kanten in einem einzigen
Energiescan abgedeckt werden, muss zusätzlich darauf geachtet werden,
die Energiebereiche sauber zu separieren, damit die Normierung einer
Kante nicht durch die nächstfolgende verfälscht wird.

\subsection{Aufbau des Startmodells}

Das Strukturmodell wird als Superzelle im POSCAR-Format vorgegeben.
Typisch ist eine $4\times4\times4$-Vervielfachung der konventionellen
Einheitszelle, was für eine FCC-Struktur $4^3 \times 4 = 256$~Atome
ergibt.
Die chemischen Elemente werden zufällig auf die Gitterplätze verteilt,
wobei die Stöchiometrie der Probe eingehalten wird.

Die Wahl der Kristallstruktur (FCC, BCC, HCP) definiert die Topologie
des Startmodells.
\evax{} kann zwar Atome verschieben, aber nicht die grundlegende
Koordination des Gitters ändern.
Daher ist es sinnvoll, mehrere Ausgangsstrukturen zu testen und
anhand der Fitresiduen zu entscheiden, welche Struktur die Daten
am besten beschreibt.

\subsection{Konfiguration und Start}

Alle Parameter werden in der Datei \texttt{parameters.dat}
zusammengefasst (vgl.\ Tabelle~\ref{tab:evax-params}).
Die Option \texttt{-autoES} veranlasst \evax{}, vor der eigentlichen
Optimierung die Werte von $S_0^2$ und $\Delta E_0$ automatisch
zu bestimmen.
Anschließend durchläuft der Algorithmus eine Screening-Phase, in der
die grobe Strukturoptimierung stattfindet, gefolgt von einer
Feinoptimierung bis zum Erreichen des \texttt{Froze\_in}-Kriteriums.

\subsection{Bewertung der Fitgüte}

Die Güte des Fits wird durch den Residuenwert $R$ quantifiziert, der
als normierte Abweichung zwischen berechnetem und experimentellem
Signal im gewählten Vergleichsraum definiert ist.
Werte unter $R \approx 0{,}1$ gelten als guter Fit; bei Multi-Edge-Fits
ist zu beachten, dass der gewichtete Gesamtresidual die
Einzelresiduen der Kanten zusammenfasst und daher höher ausfallen kann.

Neben dem numerischen Residual ist die visuelle Inspektion der
$\chi(k)$- und FT-Spektren unverzichtbar:
Systematische Abweichungen in bestimmten $k$- oder $R$-Bereichen
können auf Modelldefizite hinweisen, die im Residuenwert allein
nicht sichtbar werden.

\subsection{Interpretation der Strukturparameter}

Aus dem finalen Strukturmodell lassen sich berechnen:
\begin{itemize}[nosep]
  \item Partielle radiale Verteilungsfunktionen $g_{ij}(R)$,
  \item Koordinationszahlen $N_{ij}$ durch Integration des ersten
        Peaks von $g_{ij}(R)$,
  \item Mittlere Bindungslängen $\langle R_{ij}\rangle$ und deren
        Standardabweichungen,
  \item Warren-Cowley-Parameter $\alpha_{ij}$
        (vgl.\ Abschnitt~\ref{sec:sro}).
\end{itemize}

\begin{figure}[htbp]
  \centering
  \begin{tikzpicture}[
    node distance=1.2cm and 2cm,
    block/.style={
      draw, rounded corners, text width=5cm, align=center,
      minimum height=1cm, font=\small
    },
    arrow/.style={-{Stealth[length=6pt]}, thick}
  ]
    \node[block, fill=blue!10] (raw)
      {Rohdaten $\mu(E)$\\[0.2em]\scriptsize Athena / larch};
    \node[block, fill=blue!10, below=of raw] (chi)
      {$\chi(k)$-Extraktion\\[0.2em]\scriptsize Normierung, Hintergrund};
    \node[block, fill=green!10, below=of chi] (model)
      {Startmodell (Superzelle)\\[0.2em]\scriptsize POSCAR, $4\times4\times4$};
    \node[block, fill=orange!10, below=of model] (evax)
      {\evax{}-Optimierung\\[0.2em]\scriptsize EA + RMC, FEFF};
    \node[block, fill=red!10, below=of evax] (result)
      {Ergebnisse\\[0.2em]\scriptsize RDF, $N_{ij}$, $\alpha_{ij}$};

    \draw[arrow] (raw) -- (chi);
    \draw[arrow] (chi) -- (model);
    \draw[arrow] (model) -- (evax);
    \draw[arrow] (evax) -- (result);

    % Feedback loop
    \draw[arrow, dashed, gray] (result.east) -- ++(1.5,0)
      |- node[right, pos=0.25, font=\scriptsize\itshape, text=gray]
        {Parameter anpassen} (evax.east);
  \end{tikzpicture}
  \caption{Schematischer Workflow einer \evax{}-Analyse, von der
    Rohdatenverarbeitung bis zur strukturellen Interpretation.}
  \label{fig:workflow}
\end{figure}

% Vollständige Pipeline-Skizze
% =============================================================================
% Analyse-Pipeline: saubere vertikale Darstellung
% =============================================================================
\begin{figure}[p]
  \centering
  \begin{tikzpicture}[
    scale=0.82, transform shape,
    % Styles
    step/.style={
      draw, rounded corners=4pt, fill=#1,
      text width=12cm, align=left, minimum height=1cm,
      font=\small, inner sep=8pt
    },
    step/.default={blue!6},
    filebox/.style={
      draw, fill=yellow!12, rounded corners=2pt,
      font=\scriptsize\ttfamily, inner sep=4pt
    },
    heading/.style={
      font=\small\bfseries, anchor=west
    },
    arr/.style={-{Stealth[length=5pt]}, thick, gray!70},
    loopback/.style={-{Stealth[length=5pt]}, thick, red!50!black, rounded corners=4pt},
  ]

    % ── 1. Messung ──────────────────────────────────────────────────────────
    \node[step=blue!8] (S1) at (0, 0) {
      \textbf{1 \enspace Synchrotron-Messung}\\[3pt]
      Absorptionsspektrum $\mu(E)$ im Bereich
      \num{7500}--\SI{10000}{eV}.
      Ein einziger Energiescan deckt die K-Kanten
      von Co (\SI{7709}{eV}), Ni (\SI{8333}{eV}) und
      Cu (\SI{8979}{eV}) gleichzeitig ab.
    };

    % ── 2. Athena ───────────────────────────────────────────────────────────
    \node[step=orange!8, below=0.45cm of S1] (S2) {
      \textbf{2 \enspace Vorverarbeitung in Athena}\\[3pt]
      Import der Rohdaten, visuelle Kontrolle,
      erste $E_0$-Bestimmung.
      Export als Athena-Projekt
      \,\tikz[baseline=-0.3ex]{\node[filebox]{MEA\_athena.prj};}
    };
    \draw[arr] (S1) -- (S2);

    % ── 3. Extraktion ──────────────────────────────────────────────────────
    \node[step=orange!8, below=0.45cm of S2] (S3) {
      \textbf{3 \enspace $\chi(k)$-Extraktion mit larch}\\[3pt]
      Neu-Extraktion nötig, da Multi-Edge-Normierung in
      Athena fehlerhafte Kantensprünge lieferte.
      Pro Kante wird ein separates Energiefenster gesetzt
      (Co: ${\leq}\SI{8233}{eV}$,
       Ni: \SIrange{8100}{8879}{eV},
       Cu: ${\geq}\SI{8700}{eV}$).\\[3pt]
      Schritte: $E_0$ bestimmen $\to$
      Pre-Edge abziehen $\to$
      auf $\Delta\mu_0 = 1$ normieren $\to$
      Hintergrund entfernen (\texttt{autobk}, $R_\text{bkg}=\SI{1.0}{\AA}$)
      $\to$ Umrechnung $E \to k$.\\[4pt]
      Ausgabe:\enspace
      \tikz[baseline=-0.3ex]{\node[filebox]{MEA\_Co\_chi.dat};}\enspace
      \tikz[baseline=-0.3ex]{\node[filebox]{MEA\_Ni\_chi.dat};}\enspace
      \tikz[baseline=-0.3ex]{\node[filebox]{MEA\_Cu\_chi.dat};}\\
      {\scriptsize Jeweils zwei Spalten: $k$ (\AA$^{-1}$) und $\chi(k)$.}
    };
    \draw[arr] (S2) -- (S3);

    % ── 4. Startmodell ─────────────────────────────────────────────────────
    \node[step=green!8, below=0.45cm of S3] (S4) {
      \textbf{4 \enspace Startmodell (Superzelle)}\\[3pt]
      Kristallstruktur wählen (FCC, BCC oder HCP) und
      $4{\times}4{\times}4$-Superzelle mit 256 Atomen
      (85\,Co, 85\,Ni, 86\,Cu, zufällig verteilt) erzeugen.
      Format: POSCAR mit Gittervektoren und fraktionellen
      Koordinaten.\\[4pt]
      Ausgabe:\enspace
      \tikz[baseline=-0.3ex]{\node[filebox]{MEA\_CoNiCu.p1};}
    };
    \draw[arr] (S3) -- (S4);

    % ── 5. Konfiguration ──────────────────────────────────────────────────
    \node[step=green!8, below=0.45cm of S4] (S5) {
      \textbf{5 \enspace Konfiguration}
      \hfill\tikz[baseline=-0.3ex]{\node[filebox]{parameters.dat};}\\[3pt]
      Zentrale Steuerdatei mit allen Einstellungen:
      $k$-Bereich (\SIrange{3}{11.5}{\per\AA}),
      $k$-Gewichtung ($k^2$),
      Vergleichsraum (Wavelet),
      Streupfade ($N_\text{legs}=4$),
      FEFF-Clusterradius (\SI{6}{\AA}),
      Schrittweite, Konvergenzkriterien.\\[3pt]
      Die Option \texttt{-autoES} bestimmt $S_0^2$ und $\Delta E_0$
      automatisch vor dem Hauptlauf.
    };
    \draw[arr] (S4) -- (S5);

    % ── 6. EvAX-Kern ──────────────────────────────────────────────────────
    \node[step=red!6, below=0.45cm of S5,
          draw=red!40!black, thick] (S6) {
      \textbf{6 \enspace EvAX-Optimierung}
      \hfill{\small\texttt{./EvAX.exe parameters.dat -autoES}}\\[4pt]
      %
      \begin{minipage}[t]{5.5cm}
        \textbf{Iterationsschleife:}\\[2pt]
        {\scriptsize
        \textbf{a)} Zufälliges Atom verschieben\\
        \quad\enspace (max.\ \SI{0.005}{\AA}/Schritt)\\[1pt]
        \textbf{b)} FEFF berechnet Streuphasen\\
        \quad\enspace für die neue Umgebung\\[1pt]
        \textbf{c)} $\chi_\text{calc}(k)$ mitteln über alle\\
        \quad\enspace Absorber gleichen Typs\\[1pt]
        \textbf{d)} Vergleich mit $\chi_\text{exp}(k)$\\
        \quad\enspace im Wavelet-Raum\\[1pt]
        \textbf{e)} Metropolis-Akzeptanz:\\
        \quad\enspace besser $\to$ behalten,\\
        \quad\enspace schlechter $\to$ evtl.\ behalten\\[1pt]
        $\to$ zurück zu \textbf{a)}
        }
      \end{minipage}%
      \hfill
      \begin{minipage}[t]{5.5cm}
        \textbf{FEFF-Dateien im Arbeitsverz.:}\\[2pt]
        {\scriptsize
        \tikz[baseline=-0.3ex]{\node[filebox]{feff};} Binary (feff85L)\\[2pt]
        \tikz[baseline=-0.3ex]{\node[filebox]{run-FEFF.exe};} Shell-Wrapper\\[2pt]
        \tikz[baseline=-0.3ex]{\node[filebox]{feff\_0/ feff\_1/};} je Kante:\\
        \quad \texttt{feff.inp} -- Eingabe (Cluster)\\
        \quad \texttt{feff.bin} -- Streuphasen\\
        \quad \texttt{chi.dat} -- berechnetes $\chi(k)$\\[4pt]
        \textbf{Evolutionärer Algorithmus:}\\[2pt]
        16 parallele Strukturkopien (States).\\
        Alle $\sim$100 Iterationen: schlechte\\
        States durch gute ersetzen +\\
        leichte Mutationen.
        }
      \end{minipage}
    };
    \draw[arr] (S5) -- (S6);

    % ── 7. Ausgaben ───────────────────────────────────────────────────────
    \node[step=purple!6, below=0.45cm of S6] (S7) {
      \textbf{7 \enspace Ausgaben}
      \hfill{\scriptsize Verzeichnis \texttt{output\_files/}}\\[4pt]
      \tikz[baseline=-0.3ex]{\node[filebox]{output.dat};}
        {\scriptsize\enspace Residual vs.\ Iteration (Konvergenz)}\hfill
      \tikz[baseline=-0.3ex]{\node[filebox]{final.xyz};}
        {\scriptsize\enspace optimierte 256 Atompositionen}\\[3pt]
      \tikz[baseline=-0.3ex]{\node[filebox]{Co/EXAFS.dat};}
        {\scriptsize\enspace $\chi_\text{calc}(k)$ und $\chi_\text{exp}(k)$
         pro Kante}\hfill
      \tikz[baseline=-0.3ex]{\node[filebox]{Co/expFT.dat};}
        {\scriptsize\enspace Fouriertransformierte}\\[3pt]
      \tikz[baseline=-0.3ex]{\node[filebox]{Co/RDF\_Co\_Ni.dat};}
        {\scriptsize\enspace partielle Paarverteilungsfunktionen $g_{ij}(R)$}\hfill
      \tikz[baseline=-0.3ex]{\node[filebox]{restart*};}
        {\scriptsize\enspace Checkpoint zum Fortsetzen}
    };
    \draw[arr] (S6) -- (S7);

    % ── 8. Auswertung ────────────────────────────────────────────────────
    \node[step=green!10, below=0.45cm of S7,
          draw=green!40!black] (S8) {
      \textbf{8 \enspace Auswertung (Python)}\\[4pt]
      \begin{minipage}[t]{5.5cm}
        \tikz[baseline=-0.3ex]{\node[filebox]{plot\_results.py};}\\[2pt]
        {\scriptsize Parst \texttt{output\_files/},
         erzeugt Konvergenz-, EXAFS-,
         FT-, RDF- und SRO-Plots.}
      \end{minipage}%
      \hfill
      \begin{minipage}[t]{5.5cm}
        \tikz[baseline=-0.3ex]{\node[filebox]{parameter\_scan.py};}\\[2pt]
        {\scriptsize Automatisierte Variation aller
         Fit-Parameter (83 Kombinationen,
         6 parallele EvAX-Instanzen).}
      \end{minipage}\\[6pt]
      \textbf{Physikalische Ergebnisse:}\enspace
      {\scriptsize
      Koordinationszahlen $N_{ij}$,\enspace
      Bindungslängen $\langle R_{ij}\rangle$,\enspace
      Warren-Cowley-Parameter $\alpha_{ij}$,\enspace
      Kristallstruktur.}
    };
    \draw[arr] (S7) -- (S8);

  \end{tikzpicture}
  \caption{Vollständige Analyse-Pipeline von der Synchrotron-Messung
    bis zu den physikalischen Ergebnissen.
    Gelbe Kästen bezeichnen konkrete Dateien.
    Der rot umrandete Schritt~6 zeigt den iterativen Kern von \evax{},
    in dem FEFF-Rechnungen und RMC-Optimierung zusammenwirken.}
  \label{fig:full-pipeline}
\end{figure}


\section{Stärken, Grenzen und Fehlerquellen}

\subsection{Stärken}

\begin{itemize}[nosep]
  \item Simultane Auswertung beliebig vieler Kanten mit einem
        konsistenten Strukturmodell,
  \item Keine Parametrisierung der RDF nötig -- das Modell kann auch
        asymmetrische oder multimodale Verteilungen abbilden,
  \item Direkte Berechnung von SRO-Parametern aus der Struktur,
  \item Eingebauter Schutz vor lokalen Minima durch den EA.
\end{itemize}

\subsection{Grenzen}

\begin{itemize}[nosep]
  \item Die Rechenzeit skaliert mit der Zahl der Atome, der Kanten
        und der Berücksichtigung von MS-Pfaden; ein Lauf mit
        256~Atomen und drei Kanten benötigt typischerweise mehrere
        Stunden,
  \item Die Topologie des Gitters ist durch das Startmodell
        vorgegeben und wird nicht optimiert,
  \item Bei geringer Datenqualität (niedriges Signal-zu-Rausch-Verhältnis,
        kleiner nutzbarer $k$-Bereich) ist die Struktur
        unterbestimmt, und der Algorithmus kann verschiedene
        gleichwertige Minima finden.
\end{itemize}

\subsection{Typische Fehlerquellen}

\begin{itemize}[nosep]
  \item Fehlerhafte $\chi(k)$-Extraktion (falsche $E_0$-Zuordnung,
        unzureichende Normierung),
  \item Ungeeigneter $k$-Bereich (zu viel Rauschen bei hohem $k$,
        Hintergrundartefakte bei niedrigem $k$),
  \item Falsche Stöchiometrie oder Kristallstruktur im Startmodell,
  \item Unzureichende Konvergenz (\texttt{Froze\_in} zu klein gewählt).
\end{itemize}

% =============================================================================
% Kapitel 4: Methodik der vorliegenden Arbeit
% =============================================================================
\chapter{Methodik}
\label{ch:methodik}

\section{Proben und Messdaten}

Die Arbeit behandelt zwei Probensysteme, deren EXAFS-Daten im Rahmen
der Vorlesung bereitgestellt wurden:

\begin{description}[style=nextline]
  \item[Probe~1: \TODO{CoNiCu}-Mittelentropielegierung (MEA)]
    Eine äquimolare ternäre Legierung der 3d-Übergangsmetalle Cobalt,
    Nickel und Kupfer.
    Das Absorptionsspektrum wurde in einem einzigen Energiescan von
    ca.\ \SIrange{7500}{10000}{\electronvolt} aufgenommen und deckt
    damit die K-Kanten aller drei Elemente ab
    (Co: \SI{7709}{\electronvolt}, Ni: \SI{8333}{\electronvolt},
    Cu: \SI{8979}{\electronvolt}).
  \item[Probe~2: \TODO{CoPt}-Nanopartikel]
    Bimetallische Nanopartikel aus Cobalt und Platin.
    \TODO{Hier werden nach Erhalt der Daten die Messbedingungen,
    Kanten und Energiebereiche ergänzt.}
\end{description}

\section{Datenvorverarbeitung}
\label{sec:datenvorverarbeitung}

\subsection{$\chi(k)$-Extraktion}

Die Umwandlung der Rohdaten $\mu(E)$ in $\chi(k)$ wurde mit der
Python-Bibliothek \textsc{larch} \cite{Newville2013} durchgeführt.
Die wesentlichen Schritte sind:

\begin{enumerate}
  \item \textbf{Kantenenergie $E_0$:} Bestimmung als Wendepunkt
        der Absorptionskante.
  \item \textbf{Pre-Edge-Subtraktion:} Anpassung einer linearen
        Funktion im Bereich unterhalb der Kante und Abzug.
  \item \textbf{Normierung:} Anpassung eines Polynoms im
        Post-Edge-Bereich und Normierung auf den Kantensprung
        $\Delta\mu_0 = 1$.
  \item \textbf{Hintergrundentfernung:} Spline-Fit
        (\texttt{autobk}-Algorithmus mit $R_{\text{bkg}} = \SI{1.0}{\Angstrom}$)
        zur Bestimmung und Subtraktion des atomaren Hintergrunds
        $\mu_0(E)$.
\end{enumerate}

\subsection{Besonderheit bei Multi-Edge-Messungen}

Da bei der MEA-Probe alle drei Kanten in einem einzigen Scan vorliegen,
wurde für jede Kante ein separates Energiefenster definiert, um eine
Kontamination der Normierung durch die nächstfolgende Kante zu
vermeiden:
\begin{itemize}[nosep]
  \item Co-K: bis \SI{8233}{\electronvolt} (ca.\ \SI{100}{\electronvolt}
        vor der Ni-Kante),
  \item Ni-K: \SIrange{8100}{8879}{\electronvolt},
  \item Cu-K: ab \SI{8700}{\electronvolt}.
\end{itemize}
Die korrekte Normierung ist kritisch, da ein fehlerhafter Kantensprung
$\Delta\mu_0$ die gesamte Amplitude von $\chi(k)$ skaliert und damit
die effektiven Koordinationszahlen verfälscht.

\section{Startstrukturen}

Für die \evax{}-Analyse wurden dreidimensionale Superzellen mit
jeweils 256~Atomen erzeugt.
Die chemischen Elemente wurden stöchiometrisch korrekt und zufällig
auf die Gitterplätze verteilt.

Drei verschiedene Kristallstrukturen dienen als Startmodelle:
\begin{itemize}[nosep]
  \item \textbf{FCC} (kubisch-flächenzentriert): $4\times4\times4$
        konventionelle Einheitszellen, Gitterparameter
        $a = \SI{3.54}{\Angstrom}$, Superzellkante
        $L = \SI{14.16}{\Angstrom}$.
  \item \textbf{BCC} (kubisch-raumzentriert): $4\times4\times8$
        Einheitszellen, $a = \SI{2.87}{\Angstrom}$.
  \item \textbf{HCP} (hexagonal dichteste Packung):
        $8\times4\times4$ Einheitszellen,
        $a = \SI{2.51}{\Angstrom}$, $c = \SI{4.07}{\Angstrom}$.
\end{itemize}
Durch den Vergleich der finalen Residuen lässt sich bestimmen, welche
Struktur die experimentellen Daten am besten beschreibt.

\section{EvAX-Konfiguration}

\subsection{Fit-Parameter}

Die optimalen Fit-Parameter wurden durch eine systematische
Parameterstudie (vgl.\ Abschnitt~\ref{sec:param-scan}) ermittelt.
Tabelle~\ref{tab:fit-params} fasst die finalen Einstellungen zusammen.

\begin{table}[htbp]
  \centering
  \caption{Finale Fit-Parameter der \evax{}-Analyse.}
  \label{tab:fit-params}
  \small
  \begin{tabular}{ll}
    \toprule
    \textbf{Parameter} & \textbf{Wert} \\
    \midrule
    $k$-Bereich        & \SIrange{3}{11.5}{\per\Angstrom} \\
    $k$-Gewichtung     & $k^2$ \\
    Vergleichsraum     & Wavelet (\texttt{Space = w}) \\
    $N_{\text{legs}}$  & 4 (Mehrfachstreuung) \\
    $R_{\text{max,FEFF}}$ & \SI{6}{\Angstrom} \\
    Schrittweite       & \SI{0.005}{\Angstrom} \\
    Max.\ Verschiebung & \SI{0.4}{\Angstrom} \\
    Anzahl States      & 16 \\
    Screening-Iterationen & 2000 \\
    Froze-in           & 5000 \\
    \bottomrule
  \end{tabular}
\end{table}

\subsection{Automatische $S_0^2$- und $\Delta E_0$-Bestimmung}

Die Parameter $S_0^2$ und $\Delta E_0$ wurden über die
\texttt{-autoES}-Option von \evax{} automatisch bestimmt.
\evax{} verwendet dafür eine Vorstufe, in der diese Werte für jede
Kante separat optimiert werden, bevor die eigentliche
Strukturoptimierung beginnt.
Physikalisch sinnvolle Werte liegen bei $S_0^2 \approx 0{,}6$--$0{,}9$
und $|\Delta E_0| \lesssim \SI{15}{\electronvolt}$.

\section{Systematische Parameterstudie}
\label{sec:param-scan}

Um die Sensitivität der Ergebnisse gegenüber der Wahl der
Fit-Parameter zu untersuchen, wurde eine automatisierte
Parameterstudie durchgeführt.
In einer ersten Phase wurden die Fit-Parameter $k_{\min}$,
$k_{\max}$, $k^n$, der Vergleichsraum und $N_{\text{legs}}$
systematisch variiert (insgesamt 72~Kombinationen).
In einer zweiten Phase wurden die Optimierungsparameter
(Schrittweite, maximale Verschiebung, $R_{\max}^{\text{FEFF}}$)
mit den besten Fit-Parametern aus Phase~1 kombiniert
(11~Kombinationen).

Alle Läufe wurden parallelisiert (sechs gleichzeitige EvAX-Instanzen
in isolierten Arbeitsverzeichnissen) und mit reduzierten
Iterationszahlen (50~Screening-Iterationen, \texttt{Froze\_in}~=~50)
durchgeführt, um die Rechenzeit auf ein praktikables Maß zu begrenzen.
Die Ergebnisse dienen der Identifikation optimaler
Parametereinstellungen und der Abschätzung der Parametersensitivität,
nicht der finalen Strukturbestimmung.

\section{Strukturvergleich}
\label{sec:struct-compare}

Für die Bestimmung der zugrunde liegenden Kristallstruktur werden
Produktionsläufe mit den optimalen Parametern für jede der drei
Ausgangsstrukturen (FCC, BCC, HCP) durchgeführt.
Die Runs verwenden volle Konvergenzkriterien
(2000~Screening-Iterationen, \texttt{Froze\_in}~=~5000) und werden
anhand folgender Kriterien verglichen:

\begin{enumerate}[nosep]
  \item Gesamtresidual und kantenaufgelöste Einzelresiduen,
  \item Übereinstimmung der berechneten und experimentellen
        $\chi(k)$- und FT-Spektren,
  \item Physikalische Plausibilität der resultierenden $S_0^2$-Werte,
  \item Konsistenz der RDF-Peaks mit bekannten Bindungslängen.
\end{enumerate}

\section{Auswertung der Strukturergebnisse}

Aus den finalen \evax{}-Strukturen werden folgende Größen berechnet:

\begin{itemize}[nosep]
  \item \textbf{Partielle RDF} $g_{ij}(R)$ für alle Elementpaare,
  \item \textbf{Koordinationszahlen} $N_{ij}$ durch Integration des
        ersten Peaks,
  \item \textbf{Mittlere Bindungslängen} $\langle R_{ij}\rangle$ aus
        der Peakposition,
  \item \textbf{Warren-Cowley-Parameter} $\alpha_{ij}$
        (Gl.~\eqref{eq:warren_cowley}).
\end{itemize}

Die Unsicherheiten werden aus der Streuung der Werte über die
16~parallelen States abgeschätzt.

% =============================================================================
% Kapitel 5: Auswertung
% =============================================================================
\chapter{Ergebnisse}
\label{ch:auswertung}

% ─────────────────────────────────────────────────────────────────────────────
\section{Parameterstudie}
\label{sec:ergebnisse-paramscan}

Die systematische Variation der Fit-Parameter über 83~Kombinationen
ergab klare Trends, die in Abbildung~\ref{fig:param-scan} zusammengefasst
sind.

\begin{figure}[htbp]
  \centering
  \TODO{Hier fig8\_parameter\_scan.png und fig9\_parameter\_heatmap.png
        einfügen.}
  %\includegraphics[width=\textwidth]{fig8_parameter_scan.png}
  %\includegraphics[width=\textwidth]{fig9_parameter_heatmap.png}
  \caption{Ergebnisse der Parameterstudie.
    Oben: Residuen aller 83~Läufe, sortiert nach Gesamtresidual.
    Unten: Heatmaps der besten Residuen für Kombinationen aus
    $k$-Gewichtung und Vergleichsraum (links), $k_{\min}$ und
    $N_{\text{legs}}$ (Mitte) sowie kantenaufgelöste Residuen der
    zehn besten Läufe (rechts).}
  \label{fig:param-scan}
\end{figure}

Die wesentlichen Befunde sind:
\begin{itemize}[nosep]
  \item Der \textbf{Wavelet-Raum} liefert durchgehend niedrigere
        Residuen als der $k$- oder $R$-Raum, da er simultane Empfindlichkeit
        in beiden Dimensionen bietet.
  \item \textbf{$k^2$-Gewichtung} stellt den besten Kompromiss dar;
        $k^1$ unterschätzt die Beiträge schwerer Rückstreuer,
        $k^3$ verstärkt das Hochfrequenzrauschen.
  \item \textbf{Mehrfachstreuung} ($N_{\text{legs}} = 4$) verbessert
        den Fit gegenüber reiner Einfachstreuung ($N_{\text{legs}} = 2$)
        signifikant.
  \item Der Einfluss von $k_{\min}$ (\num{2} vs.\ \num{3}) ist gering;
        $k_{\min} = 3$ wird bevorzugt, da der Bereich $k < 3$ stärker
        durch Hintergrundartefakte kontaminiert ist.
  \item In der zweiten Phase zeigt sich, dass größere Schrittweiten
        (\SI{0.005}{\Angstrom}) und ein Clusterradius von \SI{6}{\Angstrom}
        die besten Ergebnisse liefern.
\end{itemize}

% ─────────────────────────────────────────────────────────────────────────────
\section{Probe~1: CoNiCu-Mittelentropielegierung}
\label{sec:probe1}

\subsection{Messbedingungen}

\begin{table}[htbp]
  \centering
  \caption{Messbedingungen und Datenbereich der MEA-Probe.}
  \label{tab:probe1-bedingungen}
  \small
  \begin{tabular}{lll}
    \toprule
    & \textbf{Wert} & \textbf{Anmerkung} \\
    \midrule
    Probe          & \TODO{CoNiCu MEA} & \TODO{Probenbeschreibung} \\
    Kanten         & Co-K, Ni-K, Cu-K & Simultane Messung \\
    Energiebereich & \SIrange{7500}{10000}{\electronvolt} & \\
    $k$-Bereich    & \SIrange{3}{11.5}{\per\Angstrom} & Alle drei Kanten \\
    \TODO{Temperatur} & \TODO{RT / \SI{300}{\kelvin}} & \\
    \TODO{Messgeometrie} & \TODO{Transmission / Fluoreszenz} & \\
    \bottomrule
  \end{tabular}
\end{table}

\subsection{Strukturvergleich: FCC vs.\ BCC vs.\ HCP}

Tabelle~\ref{tab:struct-comparison} vergleicht die Fitresiduen der
drei Ausgangsstrukturen.

\begin{table}[htbp]
  \centering
  \caption{Kantenaufgelöste und gewichtete Gesamtresiduen für
    verschiedene Ausgangsstrukturen der MEA-Probe.
    Der niedrigste Gesamtresidual ist hervorgehoben.}
  \label{tab:struct-comparison}
  \small
  \begin{tabular}{lccccccc}
    \toprule
    \textbf{Struktur} &
    $R_{\text{Co}}$ & $R_{\text{Ni}}$ & $R_{\text{Cu}}$ &
    $R_{\text{total}}$ &
    $S_0^2$(Co) & $S_0^2$(Ni) & $S_0^2$(Cu) \\
    \midrule
    FCC & \TODO{0.000} & \TODO{0.000} & \TODO{0.000} & \TODO{\textbf{0.000}} &
          \TODO{0.00} & \TODO{0.00} & \TODO{0.00} \\
    BCC & \TODO{0.000} & \TODO{0.000} & \TODO{0.000} & \TODO{0.000} &
          \TODO{0.00} & \TODO{0.00} & \TODO{0.00} \\
    HCP & \TODO{0.000} & \TODO{0.000} & \TODO{0.000} & \TODO{0.000} &
          \TODO{0.00} & \TODO{0.00} & \TODO{0.00} \\
    \bottomrule
  \end{tabular}
\end{table}

\TODO{Interpretation einfügen: Welche Struktur liefert den besten Fit?
Ist das Ergebnis konsistent mit der Literatur (FCC erwartet für CoNiCu)?
Wie groß ist der Unterschied?}

Abbildung~\ref{fig:struct-comparison} zeigt den direkten Vergleich der
EXAFS-Spektren und Fouriertransformierten aller drei Strukturen.

\begin{figure}[htbp]
  \centering
  \TODO{Hier fig10\_structure\_comparison.png einfügen.}
  %\includegraphics[width=\textwidth]{fig10_structure_comparison.png}
  \caption{Strukturvergleich der MEA-Probe.
    (a)~Residuen nach Kante und Struktur.
    (b)--(d)~$\chi(k)\cdot k^2$ für Co, Ni, Cu mit experimentellen
    Daten (schwarz) und Fits aus FCC (grün), BCC (blau), HCP (rot).
    (e)--(g)~Zugehörige Fouriertransformierte.
    (h)~Tabellarische Zusammenfassung der Fitparameter.}
  \label{fig:struct-comparison}
\end{figure}

\subsection{Konvergenz}

\begin{figure}[htbp]
  \centering
  \TODO{Hier Konvergenzplots einfügen (struct\_fcc\_convergence.png etc.).}
  %\includegraphics[width=0.7\textwidth]{struct_fcc_convergence.png}
  \caption{Konvergenzverlauf des Gesamtresiduals für den FCC-Fit der
    MEA-Probe.}
  \label{fig:convergence-mea}
\end{figure}

\TODO{Beschreibung des Konvergenzverhaltens: Ab welcher Iteration
stagniert der Residual? Wie vergleichen sich die drei Strukturen?}

\subsection{EXAFS-Fit und Fouriertransformation}

\begin{figure}[htbp]
  \centering
  \TODO{Hier struct\_fcc\_exafs\_ft.png einfügen.}
  %\includegraphics[width=\textwidth]{struct_fcc_exafs_ft.png}
  \caption{EXAFS-Fit der MEA-Probe (FCC-Struktur).
    Obere Reihe: $\chi(k)\cdot k^2$ (berechnet vs.\ experimentell)
    für Co, Ni und Cu.
    Untere Reihe: Zugehörige Fouriertransformierte.
    Der grau hinterlegte Bereich ($k < \SI{3}{\per\Angstrom}$) wurde
    vom Fit ausgeschlossen.}
  \label{fig:exafs-ft-mea}
\end{figure}

\TODO{Qualitative Beschreibung: Wo stimmt der Fit gut überein? Gibt es
systematische Abweichungen? Welche Kante zeigt den höchsten Residual
und warum (z.\,B.\ geringeres Signal-zu-Rausch-Verhältnis bei Cu)?}

\subsection{Radiale Verteilungsfunktionen}

\begin{figure}[htbp]
  \centering
  \TODO{Hier struct\_fcc\_rdf.png einfügen.}
  %\includegraphics[width=\textwidth]{struct_fcc_rdf.png}
  \caption{Partielle radiale Verteilungsfunktionen $g_{ij}(R)$ der
    MEA-Probe für die drei Absorbertypen.
    Jedes Panel zeigt die Beiträge der drei Nachbarelemente.}
  \label{fig:rdf-mea}
\end{figure}

\TODO{Beschreibung der RDF-Peaks: Position des ersten Peaks
(Nächste-Nachbar-Abstand), Breite (Unordnung), relative Höhe
(Koordinationszahl).}

\subsection{Koordinationszahlen und Bindungslängen}

\begin{table}[htbp]
  \centering
  \caption{Aus der \evax{}-Analyse bestimmte Koordinationszahlen und
    mittlere Bindungslängen der MEA-Probe (erste Koordinationsschale).
    \TODO{Werte nach Abschluss der Fits einfügen.}}
  \label{tab:cn-mea}
  \small
  \begin{tabular}{lSSS}
    \toprule
    \textbf{Paar} &
    {$N_{ij}$} &
    {$\langle R_{ij}\rangle$ (\si{\Angstrom})} &
    {$\sigma_{ij}$ (\si{\Angstrom})} \\
    \midrule
    Co--Co & \TODO{0.0} & \TODO{0.00} & \TODO{0.00} \\
    Co--Ni & \TODO{0.0} & \TODO{0.00} & \TODO{0.00} \\
    Co--Cu & \TODO{0.0} & \TODO{0.00} & \TODO{0.00} \\
    \midrule
    Ni--Co & \TODO{0.0} & \TODO{0.00} & \TODO{0.00} \\
    Ni--Ni & \TODO{0.0} & \TODO{0.00} & \TODO{0.00} \\
    Ni--Cu & \TODO{0.0} & \TODO{0.00} & \TODO{0.00} \\
    \midrule
    Cu--Co & \TODO{0.0} & \TODO{0.00} & \TODO{0.00} \\
    Cu--Ni & \TODO{0.0} & \TODO{0.00} & \TODO{0.00} \\
    Cu--Cu & \TODO{0.0} & \TODO{0.00} & \TODO{0.00} \\
    \bottomrule
  \end{tabular}
\end{table}

\TODO{Interpretation: Stimmen die Gesamtkoordinationszahlen
($\sum_j N_{ij} \approx 12$ für FCC)? Sind die Bindungslängen
konsistent mit dem Vegard'schen Gesetz?}

\subsection{Chemische Nahordnung}

\begin{table}[htbp]
  \centering
  \caption{Warren-Cowley-Parameter $\alpha_{ij}$ der ersten
    Koordinationsschale der MEA-Probe.
    Negative Werte zeigen Paarbildungstendenz an, positive Werte
    Vermeidung.}
  \label{tab:sro-mea}
  \small
  \begin{tabular}{lccc}
    \toprule
    & Co & Ni & Cu \\
    \midrule
    Co & \TODO{$+0.00$} & \TODO{$-0.00$} & \TODO{$\phantom{+}0.00$} \\
    Ni & \TODO{$-0.00$} & \TODO{$+0.00$} & \TODO{$\phantom{+}0.00$} \\
    Cu & \TODO{$\phantom{+}0.00$} & \TODO{$\phantom{+}0.00$} & \TODO{$+0.00$} \\
    \bottomrule
  \end{tabular}
\end{table}

\TODO{Interpretation: Welche Paare werden bevorzugt/vermieden? Ist die
Legierung ein idealer Mischkristall ($\alpha \approx 0$) oder gibt es
lokale Ordnung? Stimmen die Ergebnisse mit DFT-Vorhersagen oder
anderen experimentellen Befunden überein?}

% ─────────────────────────────────────────────────────────────────────────────
\section{Probe~2: CoPt-Nanopartikel}
\label{sec:probe2}

\subsection{Messbedingungen}

\begin{table}[htbp]
  \centering
  \caption{Messbedingungen und Datenbereich der CoPt-Probe.}
  \label{tab:probe2-bedingungen}
  \small
  \begin{tabular}{lll}
    \toprule
    & \textbf{Wert} & \textbf{Anmerkung} \\
    \midrule
    Probe          & \TODO{CoPt-Nanopartikel} & \TODO{Beschreibung} \\
    Kanten         & \TODO{Co-K, Pt-L$_3$} & \\
    Energiebereich & \TODO{...} & \\
    $k$-Bereich    & \TODO{...} & \\
    \TODO{Temperatur} & \TODO{...} & \\
    \bottomrule
  \end{tabular}
\end{table}

\subsection{Strukturvergleich}

\TODO{Analog zur MEA-Probe werden hier FCC, BCC und ggf.\ L1$_0$
verglichen. Für CoPt-Nanopartikel ist die Frage besonders relevant,
ob eine geordnete Phase (L1$_0$) vorliegt oder eine statistische
Verteilung auf dem FCC-Gitter.}

\begin{table}[htbp]
  \centering
  \caption{Strukturvergleich der CoPt-Probe.}
  \label{tab:struct-comparison-copt}
  \small
  \begin{tabular}{lcccc}
    \toprule
    \textbf{Struktur} &
    $R_{\text{Co}}$ & $R_{\text{Pt}}$ & $R_{\text{total}}$ &
    $S_0^2$(Co) \\
    \midrule
    FCC  & \TODO{0.000} & \TODO{0.000} & \TODO{0.000} & \TODO{0.00} \\
    L1$_0$ & \TODO{0.000} & \TODO{0.000} & \TODO{0.000} & \TODO{0.00} \\
    \TODO{...} & & & & \\
    \bottomrule
  \end{tabular}
\end{table}

\subsection{EXAFS-Fit und Fouriertransformation}

\begin{figure}[htbp]
  \centering
  \TODO{EXAFS- und FT-Plots der CoPt-Probe einfügen.}
  \caption{EXAFS-Fit der CoPt-Probe.
    \TODO{Analoges Layout wie Abb.~\ref{fig:exafs-ft-mea}.}}
  \label{fig:exafs-ft-copt}
\end{figure}

\subsection{Koordinationszahlen und Bindungslängen}

\begin{table}[htbp]
  \centering
  \caption{Koordinationszahlen und Bindungslängen der CoPt-Probe.}
  \label{tab:cn-copt}
  \small
  \begin{tabular}{lSSS}
    \toprule
    \textbf{Paar} &
    {$N_{ij}$} &
    {$\langle R_{ij}\rangle$ (\si{\Angstrom})} &
    {$\sigma_{ij}$ (\si{\Angstrom})} \\
    \midrule
    Co--Co & \TODO{0.0} & \TODO{0.00} & \TODO{0.00} \\
    Co--Pt & \TODO{0.0} & \TODO{0.00} & \TODO{0.00} \\
    Pt--Co & \TODO{0.0} & \TODO{0.00} & \TODO{0.00} \\
    Pt--Pt & \TODO{0.0} & \TODO{0.00} & \TODO{0.00} \\
    \bottomrule
  \end{tabular}
\end{table}

\subsection{Chemische Nahordnung}

\begin{table}[htbp]
  \centering
  \caption{Warren-Cowley-Parameter der CoPt-Probe.}
  \label{tab:sro-copt}
  \small
  \begin{tabular}{lcc}
    \toprule
    & Co & Pt \\
    \midrule
    Co & \TODO{$+0.00$} & \TODO{$-0.00$} \\
    Pt & \TODO{$-0.00$} & \TODO{$+0.00$} \\
    \bottomrule
  \end{tabular}
\end{table}

\TODO{Interpretation: Deutet der SRO-Parameter auf L1$_0$-Ordnung hin
($\alpha_{\text{Co-Pt}} < 0$, $\alpha_{\text{Co-Co}} > 0$)?
Oder handelt es sich um eine statistische Legierung?}

% ─────────────────────────────────────────────────────────────────────────────
\section{Vergleich beider Proben}
\label{sec:vergleich}

\TODO{Dieser Abschnitt soll die Ergebnisse beider Proben
gegenüberstellen. Mögliche Vergleichsaspekte:}

\begin{itemize}[nosep]
  \item Grad der chemischen Nahordnung ($\alpha$-Parameter) in beiden
        Systemen,
  \item Fitqualität in Abhängigkeit von der Anzahl der Komponenten
        (ternär vs.\ binär),
  \item Lokale Verzerrung (Streuung der Bindungslängen) als Maß für
        den Gitterverspannungseffekt,
  \item Robustheit der Ergebnisse gegenüber der Wahl der
        Ausgangsstruktur.
\end{itemize}

\begin{table}[htbp]
  \centering
  \caption{Zusammenfassender Vergleich der beiden Probensysteme.}
  \label{tab:vergleich}
  \small
  \begin{tabularx}{\textwidth}{lXX}
    \toprule
    & \textbf{CoNiCu MEA} & \textbf{CoPt NP} \\
    \midrule
    Kristallstruktur  & \TODO{FCC}  & \TODO{FCC / L1$_0$} \\
    Gesamtresidual    & \TODO{0.000} & \TODO{0.000} \\
    Nächster-Nachbar-Abstand & \TODO{\SI{0.00}{\Angstrom}} &
                               \TODO{\SI{0.00}{\Angstrom}} \\
    Gesamt-KZ (1.~Schale) & \TODO{$\approx 12$} & \TODO{$\approx 12$} \\
    Stärkster SRO-Effekt  & \TODO{Co-Ni bevorzugt} &
                             \TODO{Co-Pt bevorzugt} \\
    \bottomrule
  \end{tabularx}
\end{table}

% =============================================================================
% Kapitel 6: Diskussion
% =============================================================================
\chapter{Diskussion}
\label{ch:diskussion}

\section{Aussagekraft der Strukturbestimmung}

Die \evax{}-Analyse liefert ein atomistisches Strukturmodell, das die
experimentellen EXAFS-Spektren aller beteiligten Kanten simultan
reproduziert.
Daraus lassen sich partielle Koordinationszahlen, Bindungslängen und
die chemische Nahordnung ableiten -- Größen, die mit konventionellen
Fitmethoden bei Mehrstoffsystemen nur unter starken Annahmen
zugänglich sind.

Gleichzeitig ist zu bedenken, dass das Modell nicht eindeutig ist:
Die endliche Superzelle (256~Atome) stellt eine Stichprobe der
realen Struktur dar, und verschiedene Atomkonfigurationen können
ähnliche Residuen ergeben.
Die Verwendung des evolutionären Algorithmus mit 16~parallelen States
mildert dieses Problem, beseitigt es aber nicht grundsätzlich.
Die Streuung der Ergebnisse über die States liefert ein Maß für die
statistische Unsicherheit des Modells.

\section{Rolle der Parameteroptimierung}

Die Parameterstudie (Abschnitt~\ref{sec:ergebnisse-paramscan}) zeigt,
dass die Wahl des Vergleichsraums und die Berücksichtigung von
Mehrfachstreuung den größten Einfluss auf die Fitqualität haben.
Der Wavelet-Raum erweist sich als überlegen, was mit seiner
gleichzeitigen $k$- und $R$-Auflösung konsistent ist
(vgl.\ Abschnitt~\ref{sec:wavelet}).

Bemerkenswert ist, dass die strukturellen Ergebnisse
(Koordinationszahlen, SRO-Parameter) über die verschiedenen
Parametersätze hinweg \TODO{robust / leicht variabel} sind.
Dies spricht dafür, dass die physikalisch relevante Information
tatsächlich in den Daten enthalten ist und nicht nur ein Artefakt der
Parameterwahl darstellt.

\section{Unsicherheiten und systematische Fehler}

\subsection{$\chi(k)$-Extraktion}

Die Normierung der $\chi(k)$-Daten hat direkten Einfluss auf die
Amplitude des EXAFS-Signals und damit auf die Koordinationszahlen.
Bei der vorliegenden Multi-Edge-Messung der MEA-Probe musste die
Normierung für jede Kante separat durchgeführt werden
(Abschnitt~\ref{sec:datenvorverarbeitung}).
Die Erfahrung zeigt, dass eine um \SI{10}{\percent} fehlerhafte
Normierung zu einer entsprechenden Skalierung der
Koordinationszahlen führt.

\subsection{Korrelation von $S_0^2$ und $N$}

Der Amplitudenreduktionsfaktor $S_0^2$ und die Koordinationszahl $N$
sind physikalisch verschieden, treten in der EXAFS-Gleichung
(Gl.~\eqref{eq:exafs}) jedoch als Produkt $S_0^2 \cdot N$ auf.
Die automatische $S_0^2$-Bestimmung durch \evax{} liefert Werte, die
mit Literaturangaben für metallische Systeme konsistent sind
($S_0^2 \approx 0{,}6$--$0{,}9$).
Eine unabhängige Kalibrierung an einer Referenzprobe würde die
Zuverlässigkeit weiter erhöhen.

\subsection{Modellannahmen}

Das Strukturmodell setzt ein periodisch fortgesetztes Kristallgitter
voraus.
Für die Bulk-Legierung (Probe~1) ist diese Annahme gerechtfertigt.
Bei Nanopartikeln (Probe~2) hingegen weichen die Koordinationszahlen
an der Oberfläche deutlich von den Bulk-Werten ab, was zu einer
systematischen Unterschätzung der mittleren Koordinationszahl
führen kann.
\evax{} bietet die Möglichkeit, endliche Cluster statt periodischer
Zellen zu verwenden, was bei einer weiterführenden Analyse zu
berücksichtigen wäre.

\section{Einordnung in die Literatur}

\TODO{Nach Vorliegen der Ergebnisse: Vergleich mit veröffentlichten
Werten für CoNiCu-MEAs (z.\,B.\ Gitterparameter aus XRD, SRO aus
Monte-Carlo-Simulationen) und CoPt-Nanopartikel (L1$_0$-Ordnung,
Partikelgröße aus TEM).}

Für das CoNiCu-System ist bekannt, dass alle drei reinen Metalle
oberhalb von ca.\ \SI{700}{\kelvin} in der FCC-Phase vorliegen und
eine lückenlose Mischkristallreihe bilden \cite{George2019}.
Die Frage nach der chemischen Nahordnung -- ob bestimmte
Nachbarpaare statistisch bevorzugt oder vermieden werden -- ist
hingegen experimentell wenig untersucht und stellt daher den
eigentlichen Informationsgewinn der vorliegenden Analyse dar.

\section{Vergleich mit konventionellen Fitmethoden}

Im Vergleich zu einem konventionellen EXAFS-Fit in \textsc{Artemis}
bietet die \evax{}-Analyse zwei wesentliche Vorteile:
\begin{enumerate}[nosep]
  \item Die Zahl der effektiven Fitparameter skaliert nicht mit der
        Zahl der Streupfade, da die Atompositionen selbst die einzigen
        freien Parameter sind.
  \item Das Modell liefert unmittelbar dreidimensionale Strukturen,
        aus denen sich höhere Korrelationsfunktionen (z.\,B.\ die
        Dreikörperverteilung) extrahieren lassen.
\end{enumerate}

Der Nachteil liegt in der erheblich höheren Rechenzeit und der
Abhängigkeit vom Startmodell: Während ein Shell-Fit in Minuten
abgeschlossen ist, benötigt ein \evax{}-Lauf mit 256~Atomen und
drei Kanten mehrere Stunden.
Für Routineanalysen bleibt der konventionelle Ansatz daher das Mittel
der Wahl; für komplexe Mehrstoffsysteme, bei denen die
Parameterkorrelation zu groß wird, stellt \evax{} eine leistungsfähige
Alternative dar.

% =============================================================================
% Kapitel 7: Fazit und Ausblick
% =============================================================================
\chapter{Fazit und Ausblick}
\label{ch:fazit}

\section{Zusammenfassung}

In dieser Arbeit wurde die Software \evax{} eingesetzt, um die lokale
atomare Struktur zweier Probensysteme aus EXAFS-Daten zu bestimmen:
einer ternären CoNiCu-Mittelentropielegierung und bimetallischer
CoPt-Nanopartikel.

Die methodische Vorgehensweise umfasste die sorgfältige Extraktion
der $\chi(k)$-Daten aus Multi-Edge-Absorptionsspektren, den Aufbau
geeigneter Startstrukturen, eine systematische Parameterstudie zur
Identifikation optimaler Fit-Einstellungen sowie den Vergleich
verschiedener Kristallstrukturen (FCC, BCC, HCP).

\TODO{Nach Abschluss der Fits: Kernaussagen zu den Ergebnissen
zusammenfassen, z.\,B.:}
\begin{itemize}[nosep]
  \item \TODO{Die CoNiCu-MEA kristallisiert in der FCC-Struktur
        ($R_{\text{total}} = \ldots$ vs.\ $\ldots$ für BCC).}
  \item \TODO{Die chemische Nahordnung zeigt eine leichte
        Bevorzugung von Co-Ni-Paaren ($\alpha \approx -0{,}10$),
        während Co-Co-Paare vermieden werden
        ($\alpha \approx +0{,}11$).}
  \item \TODO{Die CoPt-Nanopartikel weisen \ldots}
\end{itemize}

Die Parameterstudie belegt, dass der Wavelet-Raum als Vergleichsraum,
eine $k^2$-Gewichtung und die Berücksichtigung von
Mehrfachstreupfaden die besten Ergebnisse liefern.
Die strukturellen Befunde erwiesen sich als robust gegenüber
Variationen im $k$-Bereich und der Schrittweite.

\section{Bewertung der Methode}

\evax{} hat sich als geeignetes Werkzeug für die Analyse komplexer
Mehrstoffsysteme erwiesen.
Die simultane Multi-Edge-Auswertung überwindet die bei konventionellen
Fits problematische Parameterkorrelation und ermöglicht die direkte
Bestimmung elementaufgelöster Strukturparameter.

Die Haupteinschränkung liegt in der Rechenzeit:
Ein einzelner Produktionslauf erfordert mehrere Stunden, und die
systematische Parameterexploration vervielfacht diesen Aufwand.
Die in dieser Arbeit entwickelte Parallelisierungsstrategie --
isolierte Arbeitsverzeichnisse mit automatisierter Konfiguration --
hat sich als praktikable Lösung bewährt und ließ sich auf einem
handelsüblichen Laptop mit sechs gleichzeitigen Prozessen umsetzen.

\section{Ausblick}

Mehrere Erweiterungen bieten sich für zukünftige Arbeiten an:

\begin{itemize}
  \item \textbf{Temperaturabhängige Messungen:}
    Durch EXAFS-Messungen bei mehreren Temperaturen ließen sich
    thermische und statische Beiträge zum Debye-Waller-Faktor
    trennen und die Debye-Temperatur elementaufgelöst bestimmen.
  \item \textbf{Finite-Size-Modelle:}
    Für die CoPt-Nanopartikel könnte ein endlicher Cluster statt
    einer periodischen Superzelle verwendet werden, um
    Oberflächeneffekte korrekt zu erfassen.
  \item \textbf{Kombination mit XRD:}
    Die Gitterkonstante aus der Röntgendiffraktometrie kann als
    zusätzlicher Constraint in die \evax{}-Optimierung einfließen
    und die Mehrdeutigkeit weiter reduzieren.
  \item \textbf{Vergleich mit DFT-Rechnungen:}
    Ab-initio-Rechnungen der Mischungsenergie und
    SRO-Gleichgewichtsparameter könnten die experimentell bestimmte
    Nahordnung unabhängig validieren.
\end{itemize}


% ── Anhang ──────────────────────────────────────────────────────────────────
\appendix
% =============================================================================
% Anhang A: Arbeitsablauf und Dateistruktur von EvAX
% =============================================================================
\chapter{Arbeitsablauf und Dateistruktur von EvAX}
\label{ch:anhang-workflow}

Dieser Anhang beschreibt die praktische Funktionsweise von \evax{}
und das Zusammenspiel mit dem ab-initio-Streucode \feff{}.
Er dokumentiert die Struktur der Ein- und Ausgabedateien und soll
den konkreten Umgang mit der Software nachvollziehbar machen.


% ─────────────────────────────────────────────────────────────────────────────
\section{Eingabedaten}

\evax{} benötigt zwei Arten von Eingaben:
experimentelle $\chi(k)$-Daten und ein atomistisches Startmodell.

\subsection{Experimentelle $\chi(k)$-Dateien}

Pro Absorptionskante wird eine zweispaltige ASCII-Datei erwartet.
Die erste Spalte enthält den Wellenzahlvektor $k$ in
\si{\per\Angstrom}, die zweite das EXAFS-Signal $\chi(k)$
(dimensionslos).
Bei einer Dreikanten-Analyse (Co, Ni, Cu) liegen also drei Dateien
vor:
\begin{center}
  \texttt{MEA\_Co\_chi.dat},\quad
  \texttt{MEA\_Ni\_chi.dat},\quad
  \texttt{MEA\_Cu\_chi.dat}
\end{center}

Wie diese Daten aus dem gemessenen Absorptionsspektrum $\mu(E)$
gewonnen werden (Normierung, Hintergrundsubtraktion, $E \to k$-Umrechnung),
ist in Abschnitt~\ref{sec:datenvorverarbeitung} beschrieben.

\subsection{Startmodell (POSCAR-Format)}
\label{sec:anhang-poscar}

Das Strukturmodell wird im POSCAR-Format angegeben, das ursprünglich
aus der Software VASP stammt.
Es definiert die Superzelle durch Gittervektoren und die Positionen
aller Atome in fraktionellen Koordinaten:

\begin{lstlisting}[caption={Aufbau einer POSCAR-Datei
  (\texttt{MEA\_CoNiCu.p1}).}, label=lst:poscar]
CoNiCu MEA fcc a=3.54        <-- Kommentarzeile
1.0                           <-- Skalierungsfaktor
  14.16  0.00   0.00          <-- Gittervektor a1 (Angstrom)
   0.00  14.16  0.00          <-- Gittervektor a2
   0.00   0.00  14.16         <-- Gittervektor a3
Co  Ni  Cu                    <-- Elementnamen
85  85  86                    <-- Anzahl pro Element (= 256)
Direct                        <-- fraktionelle Koordinaten
0.000  0.000  0.000  Cu      <-- x y z Element
0.125  0.125  0.000  Co
0.125  0.000  0.125  Co
...                           <-- 256 Zeilen insgesamt
\end{lstlisting}

Die Gittervektoren bestimmen die Geometrie der Superzelle.
Für eine FCC-Struktur mit Gitterkonstante $a = \SI{3.54}{\Angstrom}$
und $4\times4\times4$-Vervielfachung ergibt sich eine kubische Zelle
mit Kantenlänge $L = 4a = \SI{14.16}{\Angstrom}$ und
$4 \times 4^3 = 256$ Atomen.
Die Elemente werden zufällig, aber stöchiometrisch korrekt auf die
Gitterplätze verteilt.

Die Wahl der Kristallstruktur definiert die Topologie des Gitters.
\evax{} kann Atome verschieben, aber die grundlegende Koordination
nicht ändern -- daher muss die Struktur als Startannahme vorgegeben
werden.
Durch Vergleich der Fitresiduen für verschiedene Startstrukturen
(FCC, BCC, HCP) lässt sich bestimmen, welche die Daten am besten
beschreibt (Tabelle~\ref{tab:startstrukturen-anhang}).

\begin{table}[htbp]
  \centering
  \caption{Verwendete Startstrukturen mit jeweils 256 Atomen.}
  \label{tab:startstrukturen-anhang}
  \small
  \begin{tabular}{llll}
    \toprule
    \textbf{Struktur} & \textbf{Datei} &
    \textbf{Gitterkonstante} & \textbf{Superzelle} \\
    \midrule
    FCC & \texttt{MEA\_CoNiCu.p1}
        & $a = \SI{3.54}{\Angstrom}$
        & $4\times4\times4$ \\
    BCC & \texttt{MEA\_CoNiCu\_bcc.p1}
        & $a = \SI{2.87}{\Angstrom}$
        & $4\times4\times8$ \\
    HCP & \texttt{MEA\_CoNiCu\_hcp.p1}
        & $a = \SI{2.51}{\Angstrom}$, $c = \SI{4.07}{\Angstrom}$
        & $8\times4\times4$ \\
    \bottomrule
  \end{tabular}
\end{table}


% ─────────────────────────────────────────────────────────────────────────────
\section{Die Konfigurationsdatei parameters.dat}
\label{sec:anhang-params}

Alle Einstellungen eines \evax{}-Laufs werden in einer einzigen
Textdatei zusammengefasst.
Sie ist in Blöcke gegliedert, die im Folgenden erläutert werden.

\subsection{Strukturmodell}

\begin{tabular}{@{}p{5cm}p{9.5cm}@{}}
  \texttt{File\_with\_structure} &
    Pfad zur POSCAR-Datei (Startmodell). \\
  \texttt{Supercell} &
    Vervielfachung der Einheitszelle, z.\,B.\ \texttt{1x1x1} wenn
    die Superzelle bereits vollständig aufgebaut ist. \\
  \texttt{Initial\_displacement} &
    Zufällige Auslenkung aller Atome zu Beginn (\si{\Angstrom}), um
    die perfekte Gittersymmetrie zu brechen.
\end{tabular}

\subsection{Optimierung}

\begin{tabular}{@{}p{5cm}p{9.5cm}@{}}
  \texttt{Stop\_after} &
    Zahl der Iterationen in der Screening-Phase, in der die
    Streubasis häufiger aktualisiert wird. \\
  \texttt{Froze\_in} &
    Abbruchkriterium: Der Lauf endet, wenn sich der Residual über
    diese Anzahl an Iterationen nicht mehr verbessert. \\
  \texttt{Number\_of\_states} &
    Zahl der parallelen Strukturkopien im evolutionären Algorithmus
    (typisch: 16). \\
  \texttt{Maximal\_step\_length} &
    Maximale Verschiebung eines Atoms pro Iteration
    (\si{\Angstrom}). \\
  \texttt{Maximal\_displacement} &
    Maximale kumulative Verschiebung eines Atoms gegenüber der
    Startposition (\si{\Angstrom}).
\end{tabular}

\subsection{Experimentelle Daten}

\begin{tabular}{@{}p{5cm}p{9.5cm}@{}}
  \texttt{Number\_of\_spectra} &
    Anzahl der gleichzeitig gefitteten Kanten (hier: 3). \\
  \texttt{File\_with\dots EXAFS\_signal} &
    Pfade zu den $\chi(k)$-Eingabedateien. \\
  \texttt{S02} &
    Amplitudenreduktionsfaktor $S_0^2$ pro Kante.
    Wird mit der Option \texttt{-autoES} automatisch bestimmt. \\
  \texttt{dE0} &
    Energieverschiebung $\Delta E_0$ (\si{\electronvolt})
    zwischen theoretischer und experimenteller Kantenlage.
    Ebenfalls durch \texttt{-autoES} bestimmbar.
\end{tabular}

\subsection{FEFF-Konfiguration}

\begin{tabular}{@{}p{5cm}p{9.5cm}@{}}
  \texttt{FEFF\_path} &
    Pfad zum \feff{}-Binary, das die ab-initio-Streurechnung
    durchführt. \\
  \texttt{N\_legs} &
    Maximale Anzahl der Streuarme pro Pfad.
    2 = nur Einfachstreuung (Absorber $\to$ Nachbar $\to$
    zurück), 4 = einschließlich Dreifach- und
    Vierfachstreuung über Zwischenatome. \\
  \texttt{R\_max\_for\_FEFF} &
    Radius des Clusters um jeden Absorber, für den \feff{}
    Streupfade berechnet (\si{\Angstrom}).
    Größere Werte erfassen mehr Nachbarschalen, erhöhen aber
    die Rechenzeit. \\
  \texttt{Edge} &
    Absorberelemente (z.\,B.\ \texttt{Co Ni Cu}). \\
  \texttt{Edge\_type} &
    Kantentyp (z.\,B.\ \texttt{K K K}). \\
  \texttt{Update\_basis\_every} &
    Häufigkeit der Aktualisierung der FEFF-Streubasis
    (in Iterationen). \\
  \texttt{Clustering\_precision} &
    Schwelle, ab der eine veränderte lokale Umgebung als
    neuer Clustertyp erkannt wird.
\end{tabular}

\subsection{Vergleich Theorie--Experiment}

\begin{tabular}{@{}p{5cm}p{9.5cm}@{}}
  \texttt{Space} &
    Raum, in dem der Vergleich stattfindet:
    \texttt{k} ($k$-Raum), \texttt{R} ($R$-Raum) oder
    \texttt{w} (Wavelet-Raum, zweidimensional in $k$ und $R$). \\
  \texttt{k\_min}, \texttt{k\_max} &
    Grenzen des $k$-Bereichs (\si{\per\Angstrom}). \\
  \texttt{k\_power} &
    Exponent der $k$-Gewichtung ($n$ in $k^n \cdot \chi(k)$). \\
  \texttt{R\_min}, \texttt{R\_max} &
    Grenzen des $R$-Fensters für die Fouriertransformation
    (\si{\Angstrom}).
\end{tabular}


% ─────────────────────────────────────────────────────────────────────────────
\section{Ablauf eines EvAX-Laufs}
\label{sec:anhang-ablauf}

\evax{} wird von der Kommandozeile gestartet:
\begin{lstlisting}[language=bash, numbers=none]
./EvAX.exe parameters.dat -autoES
\end{lstlisting}
Die Option \texttt{-autoES} veranlasst eine Vorstufe, in der $S_0^2$
und $\Delta E_0$ für jede Kante optimiert werden, bevor die
eigentliche Strukturoptimierung beginnt.

\subsection{Initialisierung}

Die Superzelle wird aus der POSCAR-Datei gelesen und 16 Kopien
(\emph{States}) angelegt.
Die Atome werden optional um einen kleinen Betrag
(\texttt{Initial\_displacement}) zufällig ausgelenkt, um die
perfekte Gittersymmetrie zu brechen.
Anschließend führt \evax{} für jeden Absorbertyp eine erste
\feff{}-Rechnung durch, um die Streubasis aufzubauen.

\subsection{Zusammenspiel von EvAX und FEFF}

Das zentrale Element der \evax{}-Methode ist die wiederholte
Berechnung theoretischer EXAFS-Spektren aus einem sich ständig
verändernden Strukturmodell.
Dafür nutzt \evax{} den ab-initio-Code \feff{} (Version~8.5)
\cite{Rehr2010} als externen Rechenkern.
Das Zusammenspiel funktioniert wie folgt:

\begin{enumerate}
  \item \textbf{Clusterextraktion:}
    \evax{} wählt ein Absorberatom (z.\,B.\ ein Co-Atom) und
    extrahiert alle Nachbaratome innerhalb des Radius
    \texttt{R\_max\_for\_FEFF} aus der aktuellen Superzelle.
    Diese lokale Umgebung wird als Cluster in eine \feff{}-Eingabedatei
    (\texttt{feff.inp}) geschrieben.

  \item \textbf{FEFF-Rechnung:}
    \feff{} berechnet für diesen Cluster die Streuphasen und
    -amplituden aller relevanten Pfade (bis zu \texttt{N\_legs}
    Streuarme) unter Verwendung der Muffin-Tin-Approximation für
    das Kristallpotential.
    Die selbstkonsistente Berechnung der Potentiale berücksichtigt
    die tatsächliche chemische Umgebung des Absorbers.
    Das Ergebnis wird in der Binärdatei \texttt{feff.bin} abgelegt.

  \item \textbf{Spektrumberechnung:}
    \evax{} liest die Streupfade aus \texttt{feff.bin} und berechnet
    daraus das theoretische $\chi_\text{calc}(k)$ des betreffenden
    Absorberatoms.
    Über alle Absorberatome gleichen Typs wird gemittelt -- bei
    85~Co-Atomen in der Superzelle ergibt sich das mittlere
    Co-K-EXAFS als Durchschnitt über 85 individuelle Spektren.

  \item \textbf{Clusterrecycling:}
    Eine vollständige \feff{}-Rechnung für alle 256~Atome bei jeder
    Iteration wäre viel zu rechenintensiv.
    Stattdessen baut \evax{} eine Bibliothek bereits berechneter
    Cluster auf.
    Nur wenn sich die lokale Umgebung eines Absorbers über eine
    Schwelle (\texttt{Clustering\_precision}) hinaus verändert hat,
    wird eine neue \feff{}-Rechnung ausgelöst.
    Der Parameter \texttt{Update\_basis\_every} steuert zusätzlich,
    nach wie vielen Iterationen die gesamte Streubasis überprüft wird.
\end{enumerate}

\subsection{Iterationsschleife}

In jeder Iteration wählt \evax{} ein zufälliges Atom und verschiebt
es um einen zufälligen Vektor (maximale Länge:
\texttt{Maximal\_step\_length}).
Aus den neuen Atompositionen wird das theoretische EXAFS-Signal
aller Kanten berechnet und im gewählten Vergleichsraum (typischerweise
Wavelet) mit dem Experiment verglichen.
Der Residual quantifiziert die Abweichung als gewichtete Summe über
alle Kanten:
\begin{equation}
  R = \sum_{i\,\in\,\text{Kanten}} w_i \cdot
      \frac{\|\chi_\text{calc}^{(i)} - \chi_\text{exp}^{(i)}\|}
           {\|\chi_\text{exp}^{(i)}\|}
\end{equation}
Die Gewichte $w_i$ werden automatisch angepasst
(\texttt{adjust\_spectra\_weights}), damit keine Kante den Fit
dominiert.

Ob eine Verschiebung akzeptiert wird, entscheidet das
Metropolis-Kriterium:
\begin{itemize}[nosep]
  \item Sinkt der Residual, wird die Verschiebung akzeptiert.
  \item Steigt er, wird sie mit einer Wahrscheinlichkeit akzeptiert,
        die mit dem Betrag der Verschlechterung abnimmt.
        Dies verhindert, dass der Algorithmus in lokalen Minima
        stecken bleibt.
\end{itemize}

\subsection{Evolutionäre Selektion}

In regelmäßigen Abständen werden die 16~States anhand ihres Residuals
verglichen.
Strukturen mit schlechtem Fit werden durch leicht mutierte Kopien
guter Strukturen ersetzt.
Dieses Selektions-Mutations-Schema entspricht dem Prinzip der
natürlichen Auslese und erhöht die Wahrscheinlichkeit, das globale
Minimum des Residuals zu finden.

\subsection{Konvergenz und Abbruch}

Der Lauf endet, wenn der Residual sich über \texttt{Froze\_in}
aufeinanderfolgende Iterationen nicht verbessert hat.
Zusätzlich kann \texttt{Stop\_after} als Obergrenze für die
Screening-Phase gesetzt werden, nach der die Optimierung mit
seltenerer Basisaktualisierung fortfährt.


% ─────────────────────────────────────────────────────────────────────────────
\section{Dateien im Arbeitsverzeichnis}

Während eines Laufs erzeugt \evax{} eine Reihe von Dateien und
Unterordnern im Arbeitsverzeichnis.
Tabelle~\ref{tab:feff-dateien} erläutert die FEFF-bezogenen Dateien,
Tabelle~\ref{tab:evax-dateien} die sonstigen Zwischendateien.

\subsection{FEFF-bezogene Dateien}

\begin{table}[htbp]
  \centering
  \caption{Dateien und Ordner, die mit der \feff{}-Berechnung
    zusammenhängen.}
  \label{tab:feff-dateien}
  \small
  \begin{tabularx}{\textwidth}{lX}
    \toprule
    \textbf{Datei/Ordner} & \textbf{Bedeutung} \\
    \midrule
    \texttt{feff} &
      Symlink auf das \feff{}-Binary (\texttt{feff85L}).
      Berechnet Streuphasen und -amplituden ab initio unter
      Verwendung der Muffin-Tin-Approximation. \\
    \texttt{run-FEFF.exe} &
      Shell-Skript (Wrapper), das \feff{} im richtigen
      Unterverzeichnis aufruft.
      Typischer Inhalt: \texttt{cd \$1 \&\& ../feff}. \\
    \texttt{feff\_0/}, \texttt{feff\_1/}, \ldots &
      Ein Unterordner pro Absorbertyp (0~=~erste Kante, 1~=~zweite
      usw.).
      \evax{} legt hier die \feff{}-Eingabe ab und liest die
      Ergebnisse aus. \\
    \texttt{feff\_0/feff.inp} &
      Eingabedatei für \feff{}: kartesische Koordinaten des Clusters
      um einen Absorber, Muffin-Tin-Radien, Berechnungsparameter
      (Pfadfilter, $R_\text{max}$, $N_\text{legs}$). \\
    \texttt{feff\_0/feff.bin} &
      Binäre Ausgabe von \feff{}: enthält alle berechneten
      Streupfade mit Amplitude, Phasenverschiebung und effektiver
      Pfadlänge.
      Wird von \evax{} gelesen, um $\chi_\text{calc}(k)$
      zusammenzusetzen. \\
    \texttt{feff\_0/chi.dat} &
      Von \feff{} berechnetes $\chi(k)$ für einen einzelnen
      Cluster. \\
    \texttt{feff\_0/feff.log} &
      Protokoll der Rechnung (Potentialberechnung, Pfadsuche,
      Konvergenz der selbstkonsistenten Potentiale). \\
    \texttt{feff\_0/atoms.dat} &
      Atomliste des Clusters im \feff{}-internen Format. \\
    \bottomrule
  \end{tabularx}
\end{table}

\subsection{Sonstige Dateien}

\begin{table}[htbp]
  \centering
  \caption{Weitere Dateien, die \evax{} während eines Laufs erzeugt.}
  \label{tab:evax-dateien}
  \small
  \begin{tabularx}{\textwidth}{lX}
    \toprule
    \textbf{Datei} & \textbf{Bedeutung} \\
    \midrule
    \texttt{parameters.dat} &
      Zentrale Konfigurationsdatei (vgl.\
      Abschnitt~\ref{sec:anhang-params}).
      Wird bei Verwendung von \texttt{-autoES} um die bestimmten
      $S_0^2$- und $\Delta E_0$-Werte ergänzt. \\
    \texttt{window\_K\_0}, \texttt{\_1}, \texttt{\_2} &
      Hanning-Fensterfunktionen für den $k$-Raum, eine pro Kante. \\
    \texttt{window\_R\_0}, \ldots &
      Fensterfunktionen für den $R$-Raum. \\
    \texttt{window\_W\_0}, \ldots &
      Wavelet-Fensterfunktionen (deutlich größer, da
      zweidimensional in $k$ und $R$). \\
    \texttt{tempfile} &
      Temporäre Zwischendaten der aktuellen Iteration. \\
    \bottomrule
  \end{tabularx}
\end{table}


% ─────────────────────────────────────────────────────────────────────────────
\section{Ausgabedateien}
\label{sec:anhang-output}

Nach Abschluss eines Laufs enthält das Verzeichnis
\texttt{output\_files/} die Ergebnisse.
Pro Kante existiert ein Unterordner
(\texttt{Co/}, \texttt{Ni/}, \texttt{Cu/}).

\begin{table}[htbp]
  \centering
  \caption{Ausgabedateien von \evax{} und ihre Inhalte.}
  \label{tab:output-dateien}
  \small
  \begin{tabularx}{\textwidth}{lX}
    \toprule
    \textbf{Datei} & \textbf{Inhalt} \\
    \midrule
    \texttt{output.dat} &
      Konvergenzverlauf: Jede Zeile enthält die Iterationsnummer
      und den gewichteten Gesamtresidual.
      Dient zur Beurteilung, ob der Lauf konvergiert ist. \\
    \texttt{final.xyz} &
      Kartesische Koordinaten aller 256~Atome in der optimierten
      Struktur.
      Periodische Kopien (für die Visualisierung) sind angefügt;
      die ersten 256~Zeilen sind die realen Atome. \\
    \texttt{initial.xyz} &
      Koordinaten der Ausgangsstruktur vor der Optimierung. \\
    \texttt{restart\textit{N}} &
      Checkpoint nach $N$~Iterationen.
      Enthält den vollständigen Zustand (Parameter, Residuen,
      Atompositionen aller 16~States) und kann als Eingabe für
      einen Folgelauf dienen. \\
    \midrule
    \texttt{Co/EXAFS.dat} &
      Drei Spalten: $k$, $\chi_\text{calc}(k)$,
      $\chi_\text{exp}(k)$.
      Ermöglicht den direkten visuellen und quantitativen
      Vergleich. \\
    \texttt{Co/expFT.dat} &
      Fouriertransformierte: $R$, Re($\tilde\chi$),
      Im($\tilde\chi$). \\
    \texttt{Co/expBFT.dat} &
      Rücktransformierte (Back Fourier Transform) des
      gefilterten Signals. \\
    \texttt{Co/expWT.dat} &
      Wavelet-Transformierte (zweidimensionale Matrix
      in $k$ und $R$). \\
    \texttt{Co/ipolEXAFS.dat} &
      Experimentelles Signal, interpoliert auf das $k$-Gitter
      von \evax{}. \\
    \texttt{Co/RDF\_Co\_Ni.dat} &
      Partielle radiale Verteilungsfunktion $g_\text{Co-Ni}(R)$.
      Pro Absorber-Nachbar-Kombination existiert eine Datei
      (9~Dateien bei drei Elementen). \\
    \bottomrule
  \end{tabularx}
\end{table}


% ─────────────────────────────────────────────────────────────────────────────
\section{Von den Ausgabedateien zu physikalischen Größen}

Die Ausgabedateien liegen als ASCII-Text vor und können direkt
weiterverarbeitet werden.
Aus den partiellen RDF-Dateien (\texttt{RDF\_*}) lassen sich durch
Integration des ersten Peaks die Koordinationszahlen $N_{ij}$ und
aus der Peakposition die mittleren Bindungslängen
$\langle R_{ij} \rangle$ bestimmen.

Der Warren-Cowley-Parameter $\alpha_{ij}$ wird aus der
\texttt{final.xyz}-Struktur berechnet, indem für jedes Absorberatom
die Nachbarn innerhalb eines Cutoff-Radius gezählt und mit der
statistischen Erwartung verglichen werden
(vgl.\ Gl.~\eqref{eq:warren_cowley}).

Der Konvergenzverlauf (\texttt{output.dat}) zeigt, ob der
Residual ein Plateau erreicht hat -- ein Zeichen dafür, dass die
Struktur auskonvergiert ist.
Sinkt der Residual bis zum Ende noch erkennbar, sollten mehr
Iterationen (größeres \texttt{Froze\_in}) gewählt werden.

Die EXAFS- und FT-Dateien erlauben eine visuelle Kontrolle:
Systematische Abweichungen zwischen berechnetem und experimentellem
Spektrum in bestimmten $k$- oder $R$-Bereichen können auf
Modelldefizite hinweisen, die im skalaren Residualwert allein nicht
sichtbar werden.


% ── Literaturverzeichnis ────────────────────────────────────────────────────
\printbibliography[title={Literaturverzeichnis}]

% ── Abbildungsverzeichnis ───────────────────────────────────────────────────
\listoffigures

% ── Tabellenverzeichnis ─────────────────────────────────────────────────────
\listoftables

\end{document}
